\chapter{Spatial Intelligence for Embodied Systems }
\label{sec:embodied}

\FramedInitial{E}{mbodied} systems are the arena in which spatial intelligence stops being a matter of analysis and becomes a matter of consequence. An embodied agent does not merely describe a scene; it occupies one, and every perception it forms and every inference it draws is ultimately answerable to an action it must take in the physical world, under real time budgets, sensor noise, and the cost of being wrong. This chapter turns from the capabilities of the previous part to their deployment, examining how perception, reasoning, and generation are brought to bear inside agents that move, manipulate, and assist.

What makes this more than a restatement is that embodiment changes the problem. The same spatial competence looks different in a manipulator closing a grasp, a car threading traffic, a wearable inferring its user's intent, a drone holding position in wind, or a fleet of robots coordinating on a warehouse floor: each fixes its own sensors, latencies, safety margins, and notion of success. The chapter is organized by these domains rather than by capability, because the interesting content lives in the domain-specific delta, what each setting demands of spatial intelligence that a general treatment does not reveal.


\section{Robotic Manipulation}
\label{sec:p4_s9_1}



 % \RespondAuthor{ }
\subsection{Introduction}
\label{subsec:p4-sec9-1-introduction}

Robotic manipulation has advanced rapidly, driven by strides in perception, learning, and control.
Yet a significant gap remains between competent laboratory systems and robust, general-purpose robots capable of operating in unstructured settings.
A central bottleneck is that manipulation is fundamentally spatial: every decision (where to reach, how to establish contact, how to route a trajectory through clutter, or how to use a tool) hinges on perceiving, representing, and reasoning about three-dimensional structure under uncertainty.
Manipulation is thus a direct test of the \emph{spatial intelligence} developed throughout Parts~I--III: it demands that systems go beyond semantic recognition to build actionable 3D world models, ground tasks in geometric constraints, predict contact outcomes, and plan behaviors that respect the physical laws of the environment.
The perception, reasoning, and generation capabilities that such systems presuppose are surveyed in Chapter~\ref{chapter:3:sec:6}, Chapter~\ref{sec:reasoning}, and Chapter~\ref{sec:generation}; this section focuses not on re-deriving those capabilities but on how they are composed, grounded, and constrained for physical manipulation.

\subsection{System Overview}
\label{subsec:p4-sec9-1-system-overview}
A manipulation system maps raw multimodal observations and a task instruction into executable, constraint-satisfying actions, and this section organizes that mapping into four strategic \textbf{pivots}, each a fundamental shift in how systems bridge perception and physical action.
\textbf{(i) Grounding Action in Spatial Anchors} moves beyond raw pixels to vision priors and sparse spatial tokens (keypoints, dense correspondences) that stay consistent across viewpoints and instances;
\textbf{(ii) Inferring Functional Action Possibilities} shifts from object classification to extracting affordances and tool-use capabilities, bridging perception with contact-rich execution;
\textbf{(iii) Encoding Spatial Reasoning in Generalist Models} transitions from single-task policies to foundation models (VLMs/VLAs) that integrate explicit 3D representations and reasoning traces for cross-embodiment generalization; and
\textbf{(iv) Satisfying Geometric and Logical Constraints} advances from simple trajectory generation to constrained optimization and neuro-symbolic planning (TAMP), ensuring actions respect strict safety limits and long-horizon logical structure.
The four pivots are complementary and mutually reinforcing rather than strictly sequential: anchors and affordances structure the perceptual input, generalist models supply broad priors, and constraint-based planners check feasibility under the modeled geometric, physical, and logical assumptions. Figure~\ref{fig:p4_s9_1_overview} situates them as complementary capabilities within a closed-loop perception-to-action system, and the remainder of this section treats them in this order.

\begin{figure}[t]
    \centering
    \includegraphics[width=\linewidth]{figures/part4/section9/manipulation_overview.pdf}
    \caption{Spatial intelligence for robotic manipulation, illustrated by placing a mug in an open drawer: four complementary spatial capabilities link multimodal inputs to 6-DoF actions, and execution returns new observations for replanning. Evaluation bars and depth/point-cloud graphics are schematic, not measured results.}
    \label{fig:p4_s9_1_overview}
\end{figure}


\subsection{Pivotal Task-1: Grounding Action in Spatial Anchors}
\label{subsec:p4-sec9-1-pivotal-task-1-grounding-action-in-spatial-anchors}

This pivot reframes visuomotor learning from operating over raw pixels to reasoning through explicit \textbf{spatial anchors} (sparse semantic keypoints, dense correspondences, point tracks, or affordance-based contact points) that compress high-dimensional observations into spatial tokens tied to object- or task-centric reference frames and designed to remain comparatively stable across viewpoint, instance, and clutter. As Figure~\ref{fig:p4_s9_1_anchors} indicates, the approach follows a three-step logic: large-scale vision priors first \emph{produce} the anchors, policies are then \emph{conditioned} on them, and the payoff is read off as a \emph{generalization gap} between seen and unseen instances.

\begin{figure}
    \centering
    \includegraphics[width=0.98\linewidth]{figures/part4/section9/robotic_anchors.pdf}
    \caption{Grounding action in spatial anchors: (1) anchors $P_t$ produced or transferred from vision priors across views, (2) an anchor-conditioned policy, and (3) generalization measured as the seen--unseen success gap $\Delta_{\mathrm{gen}}$. Scenes, correspondences, and score bars are schematic, not measured results.}
    \label{fig:p4_s9_1_anchors}
\end{figure}

The dominant methodology leverages large-scale vision priors to generate, propagate, or retrieve these anchors, then trains standard visuomotor learners on the resulting structured inputs. A common strategy prescribes anchor points in a few reference frames and propagates them across large datasets via off-the-shelf correspondence models: P3-PO~\cite{levy2024p3poprescriptivepointpriors} propagates manually specified keypoints using dense features from DINO~\cite{dino}, DINOv2~\cite{dinov2}, or diffusion-based representations~\cite{tang2023emergentcorrespondenceimagediffusion,zhao2025high,zhang2023talefeaturesstablediffusion} (these visual representations are surveyed in Chapter~\ref{chapter:single_modality}). Relatedly, Robo-ABC~\cite{ju2024roboabcaffordancegeneralizationcategories} retrieves affordance contact points from an external memory and transfers them to novel scenes via dense semantic correspondence, while Track2Act~\cite{bharadhwaj2024track2act} learns long-term point tracks from Internet-scale video and uses them to infer object-centric transformations that serve as open-loop manipulation plans. Once anchors are obtained, policy learning conditions action prediction directly on them, through keypoint-conditioned behavior cloning (KALM~\cite{fang2024keypointabstractionusinglarge}), diffusion policies in object-relative frames (AffordDP~\cite{wu2025afforddpgeneralizablediffusionpolicy}), and shared human--robot 3D keypoint spaces (Point Policy~\cite{haldar2025pointpolicyunifyingobservations}). Across these variants, anchoring actions to spatial structure reduces perceptual ambiguity and improves robustness under distribution shift.

Because the central claim is out-of-distribution generalization, evaluation pairs high-fidelity simulation with real-robot deployment rather than relying on either alone. In simulation, affordance-centric methods such as AffordDP~\cite{wu2025afforddpgeneralizablediffusionpolicy} run on IsaacGym~\cite{makoviychuk2021isaac} over articulated objects from GAPartnet~\cite{geng2023gapartnet} whose functional parts (handles, knobs, levers) vary widely in geometry, while imitation suites such as Robomimic~\cite{robomimic2021} and RoboTwin~\cite{mu2024robotwin} isolate trajectory tracking and dual-arm coordination, with P3-PO~\cite{levy2024p3poprescriptivepointpriors} using them for controlled spatial-anchoring ablations. Real-world suites then stress in-the-wild generalization on articulated (drawers, microwaves), deformable (towel folding), and tool-use (sweeping) tasks, as in Point Policy~\cite{haldar2025pointpolicyunifyingobservations} and Track2Act~\cite{bharadhwaj2024track2act}, and a zero-shot cross-category protocol (train on one knife, test on spatulas or novel knives) is formalized by Robo-ABC~\cite{ju2024roboabcaffordancegeneralizationcategories} to measure how well retrieved anchors transfer without extra supervision. Performance is reported as Task Success Rate and, to isolate the contribution of anchoring, a Generalization Gap between seen and unseen instances. Spatially grounded policies outperform pixel baselines under shift in these studies: P3-PO reports large gains on cluttered, novel-object scenes~\cite{levy2024p3poprescriptivepointpriors}, and Robo-ABC improves both affordance-retrieval accuracy and real-robot success without heavy robot-data collection, binding demonstrations to reusable spatial structure more effectively than classic one-shot imitation~\cite{duan2017oneshotimitationlearning}.


\subsection{Pivotal Task-2: Inferring Functional Action Possibilities}
\label{subsec:p4-sec9-1-pivotal-task-2-inferring-functional-action-possibilities}

This pivot shifts from general object recognition to extracting \textbf{object affordances}, \ie, the specific action possibilities implied by an object's functional properties, and uses them as the intermediate representation that bridges raw perception and contact-rich execution. Whereas Pivotal Task-1 treats contact points as viewpoint-stable \emph{anchors} for conditioning a policy, affordances here are studied as \emph{functional possibilities} that specify what action an object supports. As Figure~\ref{fig:p4_s9_1_affordances} shows, the thread runs from how affordances are \emph{represented}, through the \emph{actions} they drive (grasping, placement, tool use), to \emph{outcome-oriented} benchmarks that score achieved function rather than imitated trajectory.

\begin{figure}
    \centering
    \includegraphics[width=0.98\linewidth]{figures/part4/section9/robotic_affordances.pdf}
    \caption{Affordances as functional action possibilities: a task- and context-conditioned affordance representation (left) supports three parallel task examples (6-DoF grasp, stable placement, hook-and-pull tool use), each judged by its functional outcome. Green checks mark intended outcomes in these illustrations, not measured results.}
    \label{fig:p4_s9_1_affordances}
\end{figure}

To capture functional properties beyond whole-object categorization, prior work builds object-centric representations (part-level segmentation~\cite{do2018affordancenet}, dense visual descriptors~\cite{florencemanuelli2018dense}, and keypoint-based encodings~\cite{manuelli2019kpam,qin2020keto,mo2021where2act}) that directly support action synthesis, and more recent vision--language models predict spatial affordance points directly from open-vocabulary prompts (RoboPoint~\cite{yuan2024robopoint}, MOKA~\cite{liu2024moka}). These representations drive learning-based 6-DoF grasping~\cite{sundermeyer2021contact,zhao2025high,jiang2021synergies,zhao2024efficient,yuan2023m2t2,deshpande2025graspmolmo} and stable object placement~\cite{zeng2020transporter,liu2022structformer,zhao2025towards,yuan2023m2t2,zhao2025anyplace} by grounding generation in functional structure rather than raw geometry, and a specialized strand further models tool--behavior relationships, either learning tool--action--effect pairs~\cite{stoytchev2008learning} or predicting specific affordances such as pulling~\cite{tikhanoff2013exploring}, so as to infer end-effector poses that achieve the intended functional outcome.

Benchmarks for affordance methods emphasize understanding object functionality over geometric matching: the property under test is \emph{functional} generalization, not trajectory imitation. For grasping, YCB-based real-world evaluations~\cite{Calli2015YCB}, GraspNet-1Billion~\cite{Fang2020GraspNet}, and simulation suites such as Dex-Net~\cite{Mahler2017DexNet} score 6-DoF stability and contact validity across novel instances and viewpoints; placement benchmarks such as RLBench~\cite{james2020rlbench} and ManiSkill~\cite{mu2021maniskill} test stable stacking and insertion in cluttered, articulated scenes, and M2T2~\cite{yuan2023m2t2} explicitly evaluates valid placement poses conditioned on object geometry and scene context. For articulated objects and tools, Where2Act~\cite{mo2021where2act} and PartNet-Mobility~\cite{xiang2020sapien} score the prediction of actionable parts and interaction parameters for pulling, pushing, and rotating, while tool-use tasks judge whether the inferred end-effector pose produces the intended causal effect (drawing, hooking, pulling) rather than reproducing a specific trajectory. Evaluation therefore centers on outcome-oriented success: stable grasps that lift and hold unseen objects under pose, scale, and appearance variation; placement success, pose accuracy, and post-release stability under clutter and occlusion; and, for articulated and tool interactions, whether the predicted parameters yield the desired functional result.



\subsection{Pivotal Task-3: Encoding Spatial Reasoning in Generalist Models}
\label{subsec:p4-sec9-1-pivotal-task-3-encoding-spatial-reasoning-in-generalist-models}

This pivot moves from specialized, single-task policies toward \textbf{generalist robotic foundation models}. Existing Vision--Language--Action (VLA) models such as RT-2~\cite{zitkovich2023rt2} and OpenVLA~\cite{kim2024openvla} show strong semantic understanding but often lack the spatial common sense needed for precise manipulation or robust generalization to novel viewpoints, so the central challenge is to strengthen spatial grounding during pre-training, enabling explicit reasoning about 3D geometry and generalization across environments and embodiments~\cite{chen2024spatialvlm}. The general foundation-model and spatial-reasoning machinery these policies inherit is treated in Chapter~\ref{sec:reasoning}; here we examine only its manipulation-specific adaptation, which Figure~\ref{fig:p4_s9_1_vla} groups into three families: explicit 3D representations, view-invariant representations, and grounded reasoning traces.

\begin{figure}
    \centering
    \includegraphics[width=0.98\linewidth]{figures/part4/section9/robotic_vla.pdf}
    \caption{Spatial grounding mechanisms for generalist robot policies: (1) explicit 3D features fused with task tokens (3D-LLM, 3D-VLA), (2) view canonicalization (OG-VLA) and spatial action discretization (SpatialVLA), two separate methods rather than stages of one, and (3) grounded reasoning traces that condition action prediction. The bottom row shows evaluation of viewpoint robustness and cross-robot transfer; all maps and overlays are schematic.}
    \label{fig:p4_s9_1_vla}
\end{figure}

The first family builds explicit 3D representations: models such as 3D-LLM~\cite{hong20233d} and 3D-VLA~\cite{zhen20243d} bridge point clouds and large language models, with 3D-VLA adding a 3D generative pre-training objective that lets the model reason about action outcomes in a consistent geometric space, effectively ``imagining'' future states. The second groups view canonicalization and spatial action encoding as distinct mechanisms: OG-VLA~\cite{singh2025og} projects multi-view observations into a point cloud and re-renders canonical orthographic views to reduce sensitivity to camera placement, while SpatialVLA~\cite{qu2025spatialvla} uses Ego3D position encoding for observations and \textbf{Adaptive Action Grids} to discretize spatial movement actions; adapting to new setups can involve re-discretizing the grids. The third injects spatial reasoning through grounded reasoning traces: ``thinking'' models such as Emma-X~\cite{sun2024emma}, MolmoAct series~\cite{lee2025molmoact, fang2026molmoact2}, closely related to TraceVLA~\cite{zheng2024tracevla}, replace direct pixel-to-action regression with intermediate, grounded outputs (2D motion trajectories or relative depth maps) that explicitly condition the final action prediction.

These approaches are evaluated under method-specific zero-shot and adaptation protocols across embodiments, viewpoints, and environments, with key metrics including multimodal goal-generation quality, robustness to extrinsic camera perturbations, and skill transfer across robot morphologies such as from a fixed-base Franka arm to a mobile manipulator~\cite{team2024octo}. The reported gains point the same way: canonical orthographic views keep grounding consistent under viewpoint change (OG-VLA), spatially aligned action grids share 3D structure across setups through discretized movement (SpatialVLA), and intermediate visual or depth-aware reasoning traces act as a regularizer on action prediction (Emma-X, MolmoAct).

\subsection{Pivotal Task-4: Satisfying Geometric and Logical Constraints}
\label{subsec:p4-sec9-1-pivotal-task-4-satisfying-geometric-and-logical-constraints}

This pivot moves planning from heuristic trajectory generation toward formal \textbf{constrained optimization} that enforces geometric feasibility, fundamental physical laws (non-penetration, friction cones, force limits), and high-level logical task structure. By casting manipulation as multi-objective optimization, these methods seek to satisfy modeled hard constraints while optimizing soft task objectives, spanning both high-frequency control and long-horizon decision-making. Figure~\ref{fig:p4_s9_1_constraints} organizes the thread along three fronts (symbolic--geometric planning, contact and trajectory optimization, and constraint-aware generative planning) before synthesizing the roles constraints play in learning-based systems.

\begin{figure}
    \centering
    \includegraphics[width=0.98\linewidth]{figures/part4/section9/robotic_constraints.pdf}
    \caption{Geometric and logical constraints in manipulation: (1) symbolic--geometric planning with revision after infeasibility, (2) contact and trajectory optimization (inset: normal force $f_n$ into the object, tangential force $f_t$ along its surface), and (3) constraint-aware generative planning. The bottom row shows three roles constraints play for learned policies.}
    \label{fig:p4_s9_1_constraints}
\end{figure}

Methodological progress has moved from purely symbolic planners toward neuro-symbolic, learning-augmented optimization frameworks that unify discrete logic with continuous physics. Classical Task and Motion Planning (TAMP) and Logic-Geometric Programming (LGP)~\cite{toussaint2015lgp} provide the foundation by jointly reasoning over symbolic task structure and geometric feasibility, and recent systems augment them with Large Language Models and Vision-Language Models for scalability in open-ended domains; for example, VIZ-COAST~\cite{yan2025vizcoast} uses the spatial reasoning of VLMs to predict refinement failures \emph{before} expensive geometric optimization, substantially reducing search complexity on long-horizon problems. At the level of physical interaction, contact-implicit optimization discovers contact schedules without explicit mode enumeration but is traditionally too slow for real-time control; GPU-parallelized systems such as SPaSM~\cite{chen2025spasm} compile constraint evaluation and gradient computation into CUDA kernels, accelerating joint optimization of object placements and robot trajectories. Beyond deterministic constraints, learning-based and generative models capture uncertainty and complex structure: differentiable physics engines~\cite{freeman2021brax} let gradients flow through physical interactions, stochastic implicit neural signed distance functions~\cite{quintero2024stochastic} model geometric uncertainty directly within the planning manifold, and diffusion-based planners such as JM2D~\cite{jung2025joint} couple model-free diffusion with model-based constraint optimization during sampling, rather than relying only on a sequential plan-then-filter interface.

Such systems are benchmarked on multi-stage re-grasping under uncertainty, forceful tool use, and tight-tolerance sequential interaction, since what they promise is reliability over long, contact-rich horizons; suites such as the Functional Manipulation Benchmark (FMB)~\cite{fmb2023} scoring not only task completion but constraint satisfaction under partial observability, external disturbance, and compounding execution error. In these systems constrained optimization plays three roles: it acts as a \textit{teacher}, generating physically valid trajectories that supervise imitation and offline reinforcement learning; as a \textit{safety shield}, where model-based projections map unconstrained neural or diffusion policy outputs onto constraint-satisfying sets at execution time~\cite{römer2025diffusionpredictivecontrolconstraints}; and as a \textit{structural prior} that prunes physically infeasible regions of the action space and curbs compounding error in long-horizon tasks. Together these roles keep learned and generative planners grounded in geometric and physical reality, reducing physically implausible outputs and supporting interpretable, feasible plans under the modeled assumptions; real-world execution still depends on state estimation, model accuracy, solver convergence, and disturbance handling.

Table~\ref{tab:p4_s9_1_benchmarks} consolidates the principal benchmarks, representative methods, and headline evaluation focus across the four pivots, situating the spatially grounded methods surveyed above within a common evaluation landscape.

\begin{table}[t]
\centering
\fontsize{8}{10}\selectfont
\setlength{\tabcolsep}{1.0mm}
\caption{Representative benchmarks, methods, and headline metrics for the four manipulation pivots, drawn from simulation and real-robot settings; metrics are reported on out-of-distribution object instances and configurations where the underlying work permits.}
\label{tab:p4_s9_1_benchmarks}
\begin{tabular}{L{0.15\textwidth} L{0.32\textwidth} L{0.25\textwidth} L{0.19\textwidth}}
\toprule
\textbf{Pivot} & \textbf{Representative Benchmark(s)} & \textbf{Representative Method(s)} & \textbf{Headline Metric} \\
\midrule
P1: Spatial anchors & IsaacGym~\cite{makoviychuk2021isaac} with GAPartNet~\cite{geng2023gapartnet}, Robomimic~\cite{robomimic2021}, RoboTwin~\cite{mu2024robotwin} + real-world & P3-PO~\cite{levy2024p3poprescriptivepointpriors}, Point Policy~\cite{haldar2025pointpolicyunifyingobservations} & Task SR / Generalization Gap \\
P2: Affordances & GraspNet-1B~\cite{Fang2020GraspNet}, Dex-Net~\cite{Mahler2017DexNet}, Where2Act~\cite{mo2021where2act}, PartNet-Mobility~\cite{xiang2020sapien} & M2T2~\cite{yuan2023m2t2} & Grasp / Placement SR \\
P3: Generalist VLA & ManiSkill~\cite{mu2021maniskill}, Octo eval~\cite{team2024octo} & SpatialVLA~\cite{qu2025spatialvla}, 3D-VLA~\cite{zhen20243d} & Zero-shot cross-embodiment SR \\
P4: Constraints & FMB~\cite{fmb2023} & SPaSM~\cite{chen2025spasm}, JM2D~\cite{jung2025joint} & Constraint satisfaction / long-horizon SR \\
\bottomrule
\end{tabular}
\end{table}

\subsection{Established Tools and Systems}
\label{subsec:p4-sec9-1-established-tools-and-systems}

Progress in spatially grounded manipulation has been inseparable from a maturing ecosystem of simulators, open software stacks, shared datasets, and, increasingly, commercial robot platforms that serve as common development and deployment substrates.

\textbf{Simulation and physics engines.} GPU-parallel engines such as NVIDIA IsaacGym and the newer Isaac Sim / Isaac Lab\footnote{\url{https://developer.nvidia.com/isaac/sim}}~\cite{makoviychuk2021isaac}, articulated-object simulators such as SAPIEN~\cite{xiang2020sapien} (with its PartNet-Mobility assets), and general-purpose physics engines including MuJoCo~\cite{todorov2012mujoco}, PyBullet\footnote{\url{https://pybullet.org}}, and the recent Genesis\footnote{\url{https://github.com/Genesis-Embodied-AI/Genesis}} provide high-throughput, contact-rich environments, while ManiSkill~\cite{mu2021maniskill}, RLBench~\cite{james2020rlbench}, Robomimic~\cite{robomimic2021}, and RoboTwin~\cite{mu2024robotwin} package standardized task suites, demonstrations, and evaluation protocols that have become de facto testbeds for imitation and policy learning.

\textbf{Open software stacks and datasets.} On the software side, the Robot Operating System (ROS\,2) with the MoveIt motion-planning framework\footnote{\url{https://moveit.ai}} remains the de facto middleware for real-robot integration, and Hugging Face's LeRobot\footnote{\url{https://github.com/huggingface/lerobot}} packages datasets, pretrained policies, and training pipelines for low-cost reproduction. Large cross-embodiment efforts such as Open X-Embodiment\footnote{\url{https://robotics-transformer-x.github.io}}, which aggregates demonstrations and the RT-X policies across dozens of robot platforms, together with open generalist checkpoints such as Octo~\cite{team2024octo} and OpenVLA~\cite{kim2024openvla} and NVIDIA's GR00T humanoid foundation models\footnote{\url{https://developer.nvidia.com/isaac/gr00t}}, expose pretrained weights and fine-tuning interfaces that lower the barrier to transferring foundation-model policies onto new morphologies.

\textbf{Commercial and humanoid platforms.} On the deployment side, a rapidly growing set of companies is moving spatially grounded manipulation out of the laboratory. Boston Dynamics\footnote{\url{https://bostondynamics.com}} (the electric Atlas and the Stretch warehouse robot), Tesla Optimus\footnote{\url{https://www.tesla.com/we-robot}}, Figure AI\footnote{\url{https://www.figure.ai}}, 1X Technologies\footnote{\url{https://www.1x.tech}}, Agility Robotics\footnote{\url{https://agilityrobotics.com}} (the Digit humanoid, piloted in logistics), Physical Intelligence\footnote{\url{https://www.physicalintelligence.company}} (the $\pi_0$ generalist policy), and Sanctuary AI\footnote{\url{https://sanctuary.ai}} each pair learned manipulation with humanoid or mobile-manipulator hardware, while firms such as Covariant\footnote{\url{https://covariant.ai}} target warehouse pick-and-place at commercial scale. These platforms are the clearest evidence that the spatial-intelligence capabilities surveyed above are crossing from benchmark tasks into economically deployed systems.

Together these tools convert the representation and reasoning machinery of Chapter~\ref{sec:reasoning} into reproducible, physically grounded manipulation pipelines; their principal limitation remains the fidelity gap between simulated and real contact dynamics, which the following discussion revisits.


\subsection{Challenges}
\label{subsec:p4-sec9-1-challenges}

Despite rapid progress across all four pivots, bringing \emph{spatial intelligence} to robust, general-purpose manipulation remains bounded by a set of intertwined difficulties that cut across perception, policy, and planning. They cluster into four threads: keeping a persistent and calibrated internal model of the scene; closing the loop between contact physics and active sensing; transferring skills across embodiments while sourcing the data to do so; and meeting the reliability budget of a deployed system.

\textbf{Persistent, calibrated world models.} The first requirement is a \emph{persistent 3D belief}: an actionable, temporally consistent world model that survives heavy occlusion, articulation, deformation, and sensor drift, and that is maintained at control rates rather than reconstructed from scratch each frame. Persistence alone is not enough if the model cannot say how much it trusts itself. \emph{Calibrated uncertainty} asks that distributional beliefs over geometry, contact state, friction, and articulation propagate through perception, policy, and planner, so that a decision reflects confidence rather than a bare point prediction and a low-confidence grasp can trigger caution instead of failure.

\textbf{Contact modeling, sensing, and active perception.} Manipulation is ultimately about contact, yet most spatial models still reason about collision-free geometry rather than the physics of touch. Progress on \emph{contact modeling and sensing} hinges on tightly integrating vision with force and tactile feedback to reason about stick--slip, compliance, and incipient failure. Because contact is often unobservable until it happens, this couples naturally to \emph{information-seeking control}: interleaving manipulation with active perception and probing (touch-to-disambiguate, look-while-acting) inside tractable belief-space planners. And because contact is precisely where simulators diverge from reality, \emph{sim-to-real for contact} stays open, demanding better material, friction, and compliance models together with fast system identification for on-robot adaptation.

\textbf{Cross-embodiment transfer and data at scale.} Skills learned on one platform rarely move cleanly to another. \emph{Cross-embodiment and cross-scene transfer} requires aligning action spaces and spatial tokens across grippers, cameras, and environments without hand-tuned frames or retraining. Sustaining that transfer is a data problem as much as an algorithmic one: \emph{data and evaluation at scale} calls for ``long-tail'' data engines that cover rare geometries, materials, and tools with counterfactual and safety labels, and for benchmarks that measure calibration, persistence, recovery, and long-horizon competence rather than success rate alone.

\textbf{Systems reliability.} Finally, a deployed manipulator must hold together as a \emph{system}: meeting latency budgets, absorbing calibration drift, detecting out-of-distribution geometry, and enabling safe fallback and continual learning without catastrophic forgetting. Closing these gaps will require tighter couplings among 3D world models, affordance predictors, language-grounded objectives, and constraint-aware planners, so that spatial representations are not only expressive but \emph{actionable, calibrated, and safe} in unstructured settings.

\newpage
\section{Autonomous Driving}
\label{sec:p4_s9_2}

% NuScenes-SpatialQA: A Spatial Understanding and Reasoning Benchmark for Vision-Language Models in Autonomous Driving
% Gpt-driver: Learning to drive with gpt


\subsection{Introduction}
\label{subsec:p4-sec9-2-introduction}
Autonomous driving sits at the intersection of spatial intelligence, transportation engineering, and human--machine interaction, and it is one of the most demanding real-world testbeds for the spatial capabilities surveyed in Part~\ref{part:3}.
A self-driving vehicle must continuously close a \emph{perceive--localize--decide--act} loop: it has to recognize and track the surrounding traffic scene, estimate its own pose within a metric map, reason about how nearby agents will move, and commit to safe control actions, all under tight latency and safety budgets.
Advances in deep learning, multi-sensor fusion, and vehicle-to-everything (V2X) communication have steadily pushed the boundaries of what vehicles can perceive, predict, and decide, while fundamental questions around robustness, trust, ethics, and regulation remain open.
Perception and reasoning themselves are treated in Chapter~\ref{chapter:3:sec:6} and Chapter~\ref{sec:reasoning}; what is specific to driving is that each capability must hold up at speed, in all weather, and with lives at stake.

Figure~\ref{fig:p4_s9_2_levels} summarizes the SAE J3016 automation levels that frame this domain.
The SAE J3016 standard \footnote{\url{https://www.sae.org/news/2019/01/sae-updates-j3016-automated-driving-graphic}} provides a widely adopted framework for defining automation levels, with most current production systems still operating at Level 2.
Progress toward higher levels, especially full Level 5 autonomy, faces major technical \cite{YurtseverLCT20}, ethical \cite{StahlTM16}, and financial \cite{fisher2020governmental} challenges.



Deep learning has driven most of this progress, from Transformers~\cite{WangLLLXLCLZH22} and graph networks~\cite{ZhangFWH22} that model multi-agent interaction to lifelong learning for open-world adaptation~\cite{ZhuBZZW24}, but its opacity raises safety concerns compounded by sensor noise and the limited reach of vehicle-to-vehicle communication, and reaching Levels 4 and 5 depends as much on redundant hardware, large-scale data pipelines, and verifiable safety assurance as on model architecture.


\begin{figure}[!ht]
    \centering
    \includegraphics[width=0.98\linewidth]{figures/part4/section9/driving_levels.pdf}
    \caption{Functional distinctions among SAE J3016 driving-automation levels. Levels 0--2 retain a human driver; Level 3 requires a human to resume driving when necessary, whereas Levels 4--5 do not rely on human fallback within their operating scope. The stair heights indicate ordered categories, not measured safety, model complexity, or technology readiness.}
    \label{fig:p4_s9_2_levels}
\end{figure}

\subsection{System Overview}
\label{subsec:p4-sec9-2-system-overview}
Modern self-driving stacks realize the \emph{perceive--localize--decide--act} loop as a modular pipeline, sketched in Figure~\ref{fig:p4_s9_2_overview}. 
A front-end \emph{perception} module converts raw, multi-modal sensor streams (camera, LiDAR, radar) into a structured model of the scene: what objects are present, where the drivable surface and lane geometry lie, and which agents are most safety-critical. 
A \emph{localization and mapping} module places the ego-vehicle within a metric, typically HD-map-anchored, representation of that scene, so that perception outputs can be expressed in a stable world frame. 
A \emph{decision and control} module then predicts how the scene will evolve and selects the steering, acceleration, and braking commands that are sent to the actuators and fed back into the environment, closing the loop. Each of the three stages is a distinct spatial-intelligence problem with its own representations, benchmarks, and failure modes, and they are the pivotal tasks below, in pipeline order.

\begin{figure}[!ht]
    \centering
    \includegraphics[width=0.98\linewidth]{figures/part4/section9/autonomous_driving_overview.pdf}
    \caption{Closed-loop spatial intelligence for autonomous driving: complementary sensor streams feed perception and localization and mapping, whose scene and ego/map state drive decision, planning, and trajectory tracking; vehicle--environment interaction returns new observations. The controller also receives ego state; the curved wire crossing is not a junction. Dashed execution logs feed evaluation separately from control. Scenes, maps, and evaluation plots are illustrative, not measured results.}
    \label{fig:p4_s9_2_overview}
\end{figure}

\subsection{Pivotal Task-1: Perception and Detection}
\label{subsec:p4-sec9-2-pivotal-task-1-perception-and-detection}
Perception is the front-end of the pipeline: it turns raw camera, LiDAR, and radar streams into a structured, layered model of the scene that every downstream module consumes, reusing the detection, segmentation, and 3D scene-understanding machinery of Chapter~\ref{chapter:3:sec:6} but specializing it to the driving Operational Design Domain (ODD). It is best read as a set of complementary structured representations, sketched in Figure~\ref{fig:p4_s9_2_perception}: object detection localizes traffic participants, semantic and occupancy prediction labels the surrounding space, lane detection recovers the road geometry that constrains motion, and pedestrian recognition isolates the most safety-critical agents.

\begin{figure}
    \centering
    \includegraphics[width=0.98\linewidth]{figures/part4/section9/driving_perception_tasks.pdf}
    \caption{Four complementary driving perception tasks feeding a shared scene model: (1) object detection, (2) semantics and 3D occupancy, (3) lane geometry, and (4) pedestrian/VRU tracking and motion prediction. Annotations and the dashed pedestrian path are illustrative, not measured predictions.}
    \label{fig:p4_s9_2_perception}
\end{figure}

The base layer is \textbf{object detection}, which identifies and localizes obstacles and is the most direct determinant of collision safety \cite{FengHWD22}. Its central practical constraint is the ODD: environmental factors such as poor weather degrade sensor visibility and model accuracy \cite{KoopmanF19}, so recent work formalizes and iteratively validates the ODD through quantifiable environmental models and scenario catalogues, trusting a detector only within conditions it has demonstrably handled, though the absence of a standardized ODD specification still makes cross-system safety claims hard to compare. Robustness to adverse weather is therefore a first-class concern \cite{zhang2023perception,TahirZACE24,HnewaR21}: rain, snow, fog, and glare corrupt LiDAR, radar, and camera signals alike, and because driving datasets under-represent these conditions, models trained on fair-weather data generalize poorly. Responses range from dedicated robust pipelines, such as multi-stage visibility enhancement coupled with deep CNN tracking \cite{HassaballahKMM21}, to new sensing hardware, \eg, frequency-modulated continuous-wave (FMCW) LiDAR and high-dynamic-range cameras that recover usable signal in glare at higher cost. Moving from the image plane into metric space, \textbf{3D object detection} localizes and tracks objects in the vehicle's surroundings \cite{LiWLXSLQD22}, where the bird's-eye-view (BEV) representation has become dominant because it yields a unified top-down map well suited to planning. BEVFormer \cite{YangCTTZZHLQLZD23} fuses multi-camera and temporal cues with spatiotemporal transformers to approach LiDAR-level accuracy on nuScenes, while Fast-BEV \cite{LiHCCLSLXSOS24} and SparseBEV \cite{LiuTLWW23} trade marginal accuracy for the lightweight, sparse, query-based computation needed for real-time on-vehicle deployment, together tracing the field's balancing of accuracy, latency, and deployability.

Detections tell the planner \emph{where objects are} but not \emph{which space is free}; that dense layer is supplied by \textbf{semantic segmentation} and its 3D extension, occupancy prediction \cite{ZhouBWN20,NestiRNBB22}. The progression runs from 2D pixel labeling with FCN~\cite{long2015fcn} and DeepLab~\cite{chen2017deeplabv2} to volumetric reasoning: the Tri-Perspective View \cite{HuangZZ0L23} lifts RGB into a 3D representation rivalling LiDAR, SurroundOcc \cite{WeiZZZZL23} uses multi-camera 2D--3D attention to capture arbitrary shapes, and occupancy world models such as OccWorld~\cite{zheng2024occworld} extend this to forecasting how the 3D scene evolves, moving beyond the limits of bounding boxes. Because these dense predictors are data-hungry and domain-sensitive, robustness techniques such as the SePiCo pixel-contrast framework \cite{XieLLLHW23} and GAN-based synthetic data \cite{AbdollahiPSMA21} are used to survive scene variability, occlusion, and weather shift. On top of the free-space layer, \textbf{lane detection} recovers the road geometry that defines the lateral reference for lane-keeping and planning; early CNN- and YOLO-style detectors trace markings well in clear conditions but are brittle on curved, worn, or occluded lanes, motivating successors that lift lanes into 3D and use attention over global context to infer occluded segments. 
Finally, \textbf{pedestrian recognition} focuses on vulnerable road users whose behavior can be difficult to predict: CNN-, YOLO-, and SSD-based detectors deliver the required precision and latency, color--thermal fusion handles low-light scenes, and spatial--temporal models couple detection with motion prediction to anticipate crossing behavior.

Across all four layers, perception is benchmarked on large-scale driving datasets, KITTI~\cite{KITTI-GeigerLU12}, the Waymo Open Dataset~\cite{WaymoOpen-SunKDCPTGZCCVHN20}, and nuScenes~\cite{caesar2020nuscenes} for 2D/3D detection, Cityscapes~\cite{cordts2016cityscapes} and BDD100K~\cite{yu2020bdd100k} for segmentation, dedicated lane benchmarks for lane detection, and CityPersons~\cite{zhang2017citypersons} for pedestrians, and scored with mean Average Precision (mAP) and the nuScenes Detection Score (NDS), mean Intersection-over-Union (mIoU), and F1/accuracy, respectively (Table~\ref{tab:p4_s9_2_benchmarks}). Increasingly these metrics are reported on adverse-weather and out-of-distribution splits, precisely to measure behavior within and beyond the declared ODD.

\subsection{Pivotal Task-2: Localization and Mapping}
\label{subsec:p4-sec9-2-pivotal-task-2-localization-and-mapping}
If perception answers \emph{what is around the vehicle}, localization and mapping answer \emph{where the vehicle is} and \emph{what the persistent structure of the world is}. The two are tightly coupled: a metric, HD-map-anchored pose lets perception outputs be expressed in a stable world frame that planning can rely on, and the map itself is built and maintained from the same sensor streams. As Figure~\ref{fig:p4_s9_2_localization} shows, the problem decomposes into pose estimation against a known map (localization), the construction of that map (3D mapping and reconstruction), and the joint solution of both when no map exists (SLAM), with multi-sensor fusion the common thread that holds accuracy together when any single modality degrades.

\begin{figure}
    \centering
    \includegraphics[width=0.98\linewidth]{figures/part4/section9/driving_localization_and_mapping.pdf}
    \caption{Localization and mapping through uncertainty-aware sensor fusion: camera, LiDAR, radar, and GNSS/IMU measurements update the state posterior, yielding a map-relative pose (evaluated by ATE/RPE), an incrementally updated 3D map, and loop closures that feed back only after verification. Pose and map are parallel products of joint estimation; only verified loop constraints return to the estimator. Recall@$N$ evaluates place retrieval, not pose accuracy.}
    \label{fig:p4_s9_2_localization}
\end{figure}

\textbf{Localization} has moved from GPS-dependent positioning to learned matching of live camera or LiDAR observations against a pre-built map, reaching centimeter-level accuracy even in GPS-denied urban canyons; semantic cues such as roads and buildings further let a vehicle re-localize in changed environments without full remapping. Because no single sensor is reliable across occlusion, weather, and GPS dropout, \textbf{sensor-fusion localization} is the dominant paradigm: probabilistic estimators (particle filters, Kalman filters, and NDT scan matching) reconcile GNSS, IMU, and LiDAR, while learned vision--radar fusion compensates for modality-specific weaknesses in rain, fog, and glare. The recurring bottlenecks here are cross-sensor spatial/temporal synchronization and fault tolerance, which motivate Bayesian uncertainty weighting and online-adaptive fusion that recalibrates as conditions change.

\textbf{3D mapping and reconstruction} supply the persistent representation that localization matches against, turning LiDAR and stereo streams into metric models; for the place-recognition and loop-closure that keep maps globally consistent, transformer descriptors such as SeqOT~\cite{ma2022seqot} match sequential LiDAR scans faster than the sensor frame rate, trained on richly annotated corpora such as ApolloScape~\cite{wang2019apolloscape}. The open problems are the cost of processing dense point clouds at city scale and the handling of occluded or incomplete returns, which favor incremental map updating over full re-mapping. When the map is unavailable, \textbf{SLAM} solves localization and mapping jointly, which is critical in GPS-denied corridors; modern driving-oriented vSLAM hybridizes classical geometry with learned features (for example, extending ORB-SLAM~\cite{mur2017orb} with learned features, semantic labels, and learned depth to favor stable landmarks) and recovers from drift through loop closure with bundle adjustment.

These capabilities are evaluated on KITTI Odometry~\cite{KITTI-GeigerLU12}, Oxford RobotCar~\cite{Maddern2017RobotCar}, and ApolloScape~\cite{wang2019apolloscape} for long-horizon SLAM and place recognition, on nuScenes~\cite{caesar2020nuscenes} and Waymo~\cite{WaymoOpen-SunKDCPTGZCCVHN20} for map-relative pose, and on Argoverse~\cite{Argoverse-ChangLSSBHW0LRH19} for the trajectory forecasting that localization feeds, \eg, with the Motion Transformer~\cite{shi2022mtr} (Table~\ref{tab:p4_s9_2_benchmarks}). The standard metrics, \ie, Absolute and Relative Pose Error (ATE/RPE), recall@$N$ for loop closure, and minADE/minFDE for forecasting, are reported under GPS-denied, adverse-weather, and sensor-degradation conditions to stress exactly the robustness the fusion machinery is built to provide.

\subsection{Pivotal Task-3: Decision Making and Control}
\label{subsec:p4-sec9-2-pivotal-task-3-decision-making-and-control}
Decision making and control is the action end of the pipeline: it consumes the scene model from perception and the pose from localization, predicts how the scene will evolve, and commits to the steering, acceleration, and braking that close the loop, drawing on the spatial-reasoning and planning paradigms of Chapter~\ref{sec:reasoning}. As Figure~\ref{fig:p4_s9_2_decision} indicates, two architectural philosophies dominate, differing in how much of the perceive-to-act mapping is learned end-to-end versus decomposed into interpretable modules, and they are increasingly converging.

\begin{figure}
    \centering
    \includegraphics[width=0.98\linewidth]{figures/part4/section9/driving_decision_and_control.pdf}
    \caption{Three alternative architectures for driving decision and control: (1) modular pipelines with an explicit BEV scene, (2) end-to-end sensor-to-control learning, and (3) hybrids with inspectable scene or trajectory intermediates and explicit constraint checks.}
    \label{fig:p4_s9_2_decision}
\end{figure}

The first philosophy is \textbf{learning-based decision making}, centered on (deep) reinforcement learning, which acquires driving policies through trial-and-error interaction and is well suited to the sequential, interactive nature of driving. 
Representative methods illustrate its design axes: the stochastic latent actor-critic (SLAC)~\cite{lee2020stochastic} learns compact latent states from high-dimensional images, PPO~\cite{schulman2017proximalpo} within a nonlinear control framework handles sharp turns and obstacle avoidance, and the DeepSet-Q architecture~\cite{huegle2019dynamic}, building on permutation-invariant set encoders~\cite{zaheer2017deep}, ingests a variable number of surrounding agents for high-level lane-change decisions.
Hybrids with explicit planning, such as Monte Carlo tree search guided by DRL or scene-representation transformers that forecast neighbor intentions, improve sample efficiency and situational awareness in dense traffic. 
The binding difficulty is that safety-critical episodes are rare in logged data, which pushes the field toward high-fidelity simulation for rare-event coverage and toward encoding hard safety constraints directly in the policy rather than trusting the reward signal alone. 
The second philosophy, \textbf{end-to-end decision making}, learns the mapping from sensors to trajectory or control directly, removing hand-designed interfaces between modules. Early instances learned state-to-action policies with DDPG~\cite{lillicrap2015continuous} in simulators such as TORCS, and multi-task backbones such as HybridNets~\cite{vu2022hybridnets} share computation across detection and segmentation without yet closing the loop to control. Current systems are trained by imitation on large driving logs: TransFuser~\cite{chitta2023transfuser} fuses camera and LiDAR features with transformer attention and predicts waypoints for a closed-loop agent, and UniAD~\cite{hu2023uniad} chains tracking, mapping, motion forecasting, and occupancy prediction through query interfaces inside one network optimized for planning. Driving world models such as GAIA-1~\cite{hu2023gaia} generate future multi-view scenes conditioned on ego actions, giving such policies a learned simulator for training and evaluation.
Opacity is the price of end-to-end training: attribution methods such as layer-wise relevance propagation (LRP)~\cite{bach2015pixel} support interpretation but do not by themselves establish safety or regulatory compliance, which is why hybrid designs keep inspectable scene or trajectory intermediates together with explicit constraint checks, whose guarantees extend only as far as the modeled constraints.

Closed-loop driving policies are evaluated chiefly in simulation, \ie, the CARLA urban simulator~\cite{dosovitskiy2017carla} (NoCrash, Longest6), the nuPlan planning benchmark~\cite{caesar2021nuplan}, and highway-interaction simulators built on HighD/INTERACTION (Table~\ref{tab:p4_s9_2_benchmarks}), because rare, safety-critical events cannot be exercised safely on public roads. The headline metrics are the CARLA Driving Score (route completion penalized by infractions), collision and intervention rates, and route completion, complemented by comfort proxies such as jerk and, for the open-loop forecasting sub-problem, minADE/minFDE.

\begin{table}[!th]
\centering
\fontsize{8}{10}\selectfont
\setlength{\tabcolsep}{1.0mm}
\caption{Representative benchmarks for spatial intelligence in autonomous driving, covering perception, localization and mapping, motion prediction, and decision making.}
\label{tab:p4_s9_2_benchmarks}
\begin{tabular}{L{0.22\textwidth} L{0.22\textwidth} L{0.28\textwidth} L{0.19\textwidth}}
\toprule
\textbf{Pivotal task} & \textbf{Benchmark} & \textbf{Representative method} & \textbf{Headline metric} \\
\midrule
Perception (3D detection) & nuScenes & BEVFormer~\cite{YangCTTZZHLQLZD23} & NDS / mAP \\
Perception (3D detection) & Waymo Open & SparseBEV~\cite{LiuTLWW23} & mAP / latency \\
Perception (segmentation) & Cityscapes & SePiCo~\cite{XieLLLHW23} & mIoU \\
Perception (occupancy) & nuScenes Occ3D & SurroundOcc~\cite{WeiZZZZL23} & mIoU \\
Localization \& mapping & KITTI Odometry & Deep / semantic vSLAM & ATE / RPE \\
Prediction (forecasting) & Argoverse & Motion Transformer (MTR)~\cite{shi2022mtr} & minADE / minFDE \\
Decision \& control & CARLA (Longest6) & TransFuser~\cite{chitta2023transfuser} & Driving Score \\
Decision \& control & nuScenes (open-loop planning) & UniAD~\cite{hu2023uniad} & L2 error / collision rate \\
Decision \& control & nuPlan & Learned planners & Closed-loop score \\
\bottomrule
\end{tabular}
\end{table}


\subsection{Established Tools and Systems}
\label{subsec:p4-sec9-2-established-tools-and-systems}
The systems on the road today populate a deployment ladder whose rungs are the SAE levels of Figure~\ref{fig:p4_s9_2_levels}, but what separates one rung from the next is rarely a single breakthrough in perception accuracy.
It is, instead, two coupled engineering choices: how wide an Operational Design Domain (ODD) the system claims, and where legal responsibility for the driving task sits.
Three tiers are in genuine production, \ie, supervised Level~2 assistance, certified Level~3 ``eyes-off'' automation, and driverless Level~4 robotaxis, and, as the following progression shows, moving up the ladder consistently means adding sensing and compute redundancy while \emph{narrowing} the domain in which the autonomy guarantee holds.

\subsubsection{Large-scale Production ADAS (SAE Level 2 / 2+)}
The defining property of Level~2 is that the system performs continuous lateral and longitudinal control, yet the human remains the legal driver and the immediate fallback; consequently, every Level~2 product is fundamentally an exercise in bounding the ODD and enforcing driver attention rather than in achieving autonomy, and the deployed systems differ mainly in how they strike that balance.
One design philosophy gates hands-free operation to pre-validated, HD-mapped limited-access highways and polices it with in-cabin driver monitoring: \textbf{GM Super Cruise}\footnote{\url{https://www.gm.com/innovation/autonomous-driving}}, \textbf{Ford/Lincoln BlueCruise}\footnote{\url{https://www.lincoln.com/technology/bluecruise/}}, and \textbf{BMW Driving Assistant Professional}\footnote{\url{https://www.bmwgroup.com/en/news/general/2023/driver-assistance-systems.html}} all keep the vehicle centered, regulate following distance, and offer assisted lane changes, but escalate alerts and disengage the moment attention drifts or the car leaves a supported segment; their reliability is a direct product of this deliberately narrow envelope.
A second philosophy trades map dependence for scale: \textbf{Tesla Autopilot}\footnote{\url{https://www.tesla.com/en_gb/support/autopilot}} and \textbf{FSD (Supervised)}\footnote{\url{https://www.tesla.com/en_gb/fsd}} rely on a camera-first vision stack refined through fleet-wide over-the-air training, which extends assistance to a far broader and less-bounded set of scenarios but, lacking comparable geofencing, still demands constant supervision and degrades on poor markings, construction zones, and adverse weather. At the conservative end, \textbf{Toyota Safety Sense} declines hands-free operation altogether, functioning as an always-on crash-avoidance layer (automatic emergency braking, lane-keeping, and adaptive cruise) that intervenes only to prevent or mitigate collisions. Despite these differing emphases, the binding limitation is identical across all of them: a bounded ODD coupled with a mandatory human fallback, and it is precisely that fallback which the next tier begins to remove.

\subsubsection{Certified ``Eyes-off'' Conditional Automation (SAE Level 3)}
Level~3 is the first tier at which, inside a defined envelope, the driver may legally divert attention and the manufacturer assumes liability for the driving task. This shift in responsibility, rather than any new perception capability, is what dictates the engineering.
\textbf{Mercedes-Benz Drive Pilot}\footnote{\url{https://www.mercedes-benz.co.uk/passengercars/technology/autonomous-driving.html}} and \textbf{Honda Traffic Jam Pilot}\footnote{\url{https://global.honda/en/tech/Traffic_Jam_Pilot_TJP/}} purchase the ``eyes-off'' guarantee by pairing complementary sensing (cameras, radar, and LiDAR with high-precision, HD-map-anchored localization) with redundant braking, steering, and power paths that retain control through a single-component failure.
Both, tellingly, keep the envelope deliberately small (low-speed or traffic-jam flow on mapped, divided highways), and both center on the same safety mechanism: as conditions approach the ODD boundary, the system issues a timed takeover request and, absent a driver response, executes a minimal-risk maneuver such as a controlled stop with hazard signals.
The narrowness of the Level~3 envelope is therefore not a marketing concession but the direct cost of guaranteeing safe behavior once the human is no longer continuously in the loop: a cost the next tier pays in full by removing the driver entirely.

\subsubsection{Driverless Robotaxis and Autonomous Ride-hailing (SAE Level 4)}
Level~4 removes the human from the loop entirely within a geofence: there is no driver to hand back to, so the vehicle must resolve every in-domain situation itself or bring itself to a safe state, with remote operators supplying only high-level guidance rather than moment-to-moment control.
To meet this far stricter bar, \textbf{Waymo One}\footnote{\url{https://waymo.com/}} and \textbf{Baidu Apollo Go}\footnote{\url{https://www.apollogo.com/}} converge on essentially the same recipe: a dense, overlapping multi-sensor suite (lidar, radar, and camera vision) for redundant coverage, HD-map-anchored localization, and conservative prediction-and-planning that favors lawful, cautious maneuvers through dense urban traffic.
Each makes the problem tractable not by claiming universal competence but by tightly restricting where and when it operates and validating performance city by city before extending the geofence.
Across the three tiers, every step up the ladder buys wider autonomy by narrowing the operating domain and adding redundancy, until at Level~4 the system finally internalizes the fallback that Levels~2 and~3 delegated to the driver.



\subsection{Challenges}
\label{subsec:p4-sec9-2-challenges}
Every tier of the deployment ladder works because it limits where it operates, and general, trustworthy autonomy requires lifting those limits. Four threads, each cutting across perception, localization, and decision, stand in the way: justifying life-critical decisions (\emph{safety and explainability}), dependable behavior across the long tail of road, weather, and traffic conditions (\emph{resilience and robustness}), adaptation to genuinely novel situations from little prior data (\emph{unfamiliar road scenarios}), and continual learning on deployed vehicles without eroding mastered behavior (\emph{enhanced autonomy in complex scenarios}).

\paragraph{Safety and Explainability.}
Because driving decisions are life-critical, the inability to explain why a model selected a given maneuver is a first-order obstacle to certification and public trust.
Post-hoc attribution methods such as layer-wise relevance propagation~\cite{bach2015pixel} and Shapley-value attribution expose which inputs drove a decision, supporting both engineering debugging and regulatory audit, while causal-inference models go further by separating mere correlation from the actual influence of factors such as road state and a maneuver's effect on surrounding agents.
The unresolved difficulty is that these explanations remain largely retrospective and local to individual decisions; converting them into guarantees that provably hold across the full operating domain is still an open problem.

\paragraph{Resilience and Robustness.}
Robustness to the long tail of driving conditions is commonly pursued through curriculum learning, where a policy is trained on a graded sequence from simple navigation to dense mixed-traffic interaction, which improves sample efficiency and training stability relative to confronting hard cases directly.
Its effectiveness, however, hinges on a scenario curriculum that faithfully spans the deployment distribution and on objectives that keep safety and comfort constraints satisfied at every difficulty level.
The practical risk is that, absent a benchmark which deliberately surfaces rare but critical events, strong curriculum performance on common scenes can mask failures precisely on the cases that determine real-world safety.

\paragraph{Unfamiliar Road Scenarios.}
A distinct difficulty is generalization to genuinely novel situations (an unseen sign, an irregular work zone, an atypical vehicle) where data-hungry models degrade but human drivers adapt from a single exposure.
Approaches that target this gap include probabilistic program induction in a Bayesian framework, which represents a new concept as a compact program inferred from few examples, and the injection of compositional, causal, and meta-learning priors so that, for instance, a system can infer that a fallen obstacle will trigger sudden lane changes in adjacent traffic.
The standing obstacle is validating such few-shot behavior to a safety standard, since the very rarity of these events makes their failure modes hard to measure and certify.

\paragraph{Enhanced Autonomy in Complex Scenarios.}
Sustained operation in changing environments motivates lifelong learning~\cite{ZhuBZZW24}, in which a deployed system keeps acquiring new knowledge (a new traffic pattern or road layout) without retraining from scratch and without overwriting prior competence.
The central technical barrier is catastrophic forgetting, which is addressed by memory-augmented architectures and continual-learning regularizers that protect previously learned parameters, together with transfer mechanisms that reuse knowledge across tasks and cities.
The harder, still-open requirement is to perform such updates on-vehicle under real-time and verification constraints, so that an update improving one scenario cannot silently regress a safety-critical behavior elsewhere.

\newpage

\section{Wearable AI System}
\label{sec:p4_s9_3}


\subsection{Introduction}
\label{subsec:p4-sec9-3-introduction}

Wearable devices have become a central sensing modality for embodied and spatial intelligence, delivering continuous, fine-grained measurements of human motion, physiology, and physical interaction. Because they are worn on or integrated into the body, they access proprioceptive and kinesthetic cues, joint rotations, muscle activations, contact forces, and locomotion patterns, which exteroceptive sensors such as RGB-D cameras or LiDAR cannot capture directly. We use the term broadly, covering inertial measurement units (IMUs), surface electromyography (EMG) sensors, smart textiles, haptic interfaces, and wearable cameras: body-mounted technologies that together model human dynamics and embodied interaction in real-world settings. It is this egocentric, body-centric vantage point that makes wearables valuable for spatial intelligence, letting a system reason about how a person moves through, reads, and acts upon their surroundings, and it has driven rapid adoption across health monitoring, sports analytics, rehabilitation~\cite{patel2012review}, industrial safety, and human--robot collaboration.

The signals wearables provide are especially valuable because they are anticipatory. Joint kinematics and muscle activity reveal what a person is doing and, crucially, what they are about to do next, which is what underpins intent recognition, action forecasting, and safe coordination in shared spaces. Wearable IMU arrays, for example, recover whole-body skeletal motion in cluttered or occluded settings where optical capture breaks down~\cite{huang2018dip, vonmarcard2017sparse}, and EMG surfaces pre-activation patterns that precede any visible movement~\cite{atzori2014ninapro}.

These devices are often bracketed with head-mounted XR hardware, yet their role in spatial intelligence differs in kind. XR headsets integrate sensing and display for immersive visualization and mixed-reality interaction, leaning on outward-facing cameras for SLAM and environment mapping. The wearables treated here look inward instead, supplying the biomechanical and physiological state, high-frequency, continuous, and on-body, that an XR headset does not observe.

The spatial-intelligence capabilities that wearables ultimately serve, pose estimation, action recognition, localization, and trajectory prediction, are formulated in Part~III (spatial perception in Chapter~\ref{chapter:3:sec:6}, spatial reasoning in Chapter~\ref{sec:reasoning}). This section does not re-derive them; it concentrates on the body-centric sensing layer, the devices that acquire embodied signals, the systems assembled around them, and the tasks, methods, and failure modes that surface once spatial intelligence must be grounded in continuous on-body measurement, spanning long-horizon activity and gait analysis, intention prediction for human--robot collaboration, and real-time safety in industrial and logistics settings.


\subsection{Wearable Devices}
\label{subsec:p4-sec9-3-wearable-devices}

Body-centric sensing spans a diverse device ecosystem, summarized in Table~\ref{tab:p4_s9_3_devices}.

\paragraph{Motion and Body-State Sensing.}
Inertial measurement units (accelerometers, gyroscopes, and often magnetometers) are the most widely deployed body-worn sensors. Consumer watches and bands such as the Apple Watch\footnote{\url{https://www.apple.com/watch/}}, Samsung Galaxy Watch\footnote{\url{https://www.samsung.com/global/galaxy/galaxy-watch/}}, Fitbit\footnote{\url{https://www.fitbit.com/}}, and Xiaomi Band stream inertial data for activity and gait monitoring, while research and industrial systems place IMUs at specific anatomical sites and fuse them into full-body kinematics~\cite{vonmarcard2017sparse, yi2021transpose}; because they need no line of sight, they keep working under occlusion and poor lighting. Full-body motion-capture suits such as Xsens\footnote{\url{https://www.movella.com/products/xsens}}, Rokoko Smartsuit Pro\footnote{\url{https://www.rokoko.com/}}, and Noitom\footnote{\url{https://www.noitom.com/}} combine many IMUs, sometimes with pressure or stretch sensors, into accurate skeletal pose outside the studio, which has made them standard in biomechanics, sports, robotics, and character animation~\cite{trumble2017totalcapture}. Smart clothing embeds conductive fabrics, stretch sensors, or distributed pressure sensors in garments for unobtrusive, long-duration monitoring of posture, joint angles, and whole-body contact~\cite{luo2021learning}, and physiological sensors (heart rate, galvanic skin response, respiration), foot-mounted inertial or pressure sensors, and ear-worn IMUs add workload, ground-contact timing, and head orientation~\cite{patel2012review}.

\paragraph{Neuromuscular, Egocentric, and Interactive Wearables.}
Electromyography (EMG) armbands and high-density arrays sense the muscle activations that precede visible movement, giving early cues for gesture recognition and intent estimation~\cite{atzori2014ninapro}, and combining EMG with inertial data makes them more robust to arm posture and fatigue~\cite{atzori2016deep}. Wearable cameras add the exteroceptive, first-person view: egocentric RGB, depth, or fisheye streams supply scene geometry and activity context that can be fused with on-body signals~\cite{grauman2022ego4d, adt}. Two device families also act on the body. Haptic gloves such as HaptX\footnote{\url{https://haptx.com/}} and SenseGlove\footnote{\url{https://www.senseglove.com/}} and haptic vests such as bHaptics\footnote{\url{https://www.bhaptics.com/}} return tactile, vibration, or force feedback during manipulation, teleoperation, and training~\cite{pacchierotti2017wearable}, and exoskeletons combine IMUs, torque and force sensors, and joint encoders with actuation, estimating user intent in real time to assist or rehabilitate movement~\cite{young2017state}, which makes them a direct test of embodied spatial intelligence.

\begin{table}[!th]
\centering
\caption{Wearable body-centric sensing devices, by sensed signal and recovered spatial or biomechanical quantity. Haptic and exoskeletal wearables are partly \emph{actuators} that return force and tactile signals to the wearer.}
\label{tab:p4_s9_3_devices}
\fontsize{8}{10}\selectfont
\setlength{\tabcolsep}{1.0mm}
\begin{tabular}{L{0.22\textwidth} L{0.28\textwidth} L{0.42\textwidth}}
\toprule
\textbf{Device family} & \textbf{Primary signal} & \textbf{Spatial / biomechanical quantity recovered} \\
\midrule
IMU wearables            & accel., gyro., (magnet.)          & orientation, acceleration, gait, sparse-to-full body kinematics \\
Motion-capture suits     & multi-IMU (+ pressure/stretch)    & full-body skeletal pose in unconstrained settings \\
EMG / gesture bands      & neuromuscular activation          & muscle intent, hand/arm gestures, pre-movement cues \\
Smart clothing / textile & distributed pressure, strain      & posture, joint angles, whole-body contact patterns \\
Egocentric cameras       & first-person RGB / depth          & scene geometry, activity, environmental context \\
Haptic / force-feedback  & tactile, vibration, force (I/O)   & contact rendering, spatial alerts (feedback channel) \\
Exoskeletons             & IMU, torque, force, joint enc.\   & joint torque, user effort and movement intent \\
Physiological / other    & HR, GSR, respiration; foot press. & workload, exertion, ground contact, head orientation \\
\bottomrule
\end{tabular}
\end{table}


\subsection{System Overview}
\label{subsec:p4-sec9-3-system-overview}

Taken together, the devices above form the front end of a single body-centric perception-to-action loop that defines the wearable AI system as a whole. 
Raw on-body signals are first \emph{sensed} (inertial, neuromuscular, tactile, and egocentric-visual streams), then \emph{modeled} into geometric and biomechanical state (joint kinematics, body pose, contact forces), \emph{interpreted and predicted} into higher-level behavioral and intentional descriptions (activities, gestures, short-horizon motion forecasts), and finally fed back to the wearer or to a collaborating agent through \emph{action} channels (haptic cues, exoskeleton assistance, safety alerts, or shared-control commands). 
The loop then closes: the consequences of feedback alter the wearer's motion, which the sensors observe anew.

Figure~\ref{fig:p4_s9_3_overview} schematizes this loop and situates the three pivotal tasks that organize the remainder of this section as complementary stages of it. 
\textbf{Pivotal Task-1 (Embodied Motion Sensing and Pose Reconstruction)} occupies the sense-to-model stage, turning body-centric signals into stable geometric descriptions of motion. 
\textbf{Pivotal Task-2 (Behavioral Understanding and Intention Modeling)} occupies the interpret-and-predict stage, mapping reconstructed motion and neuromuscular cues to recognized activities and anticipated actions. 
\textbf{Pivotal Task-3 (Safe Collaboration and Embodied Interaction)} closes the loop at the action stage, converting interpreted human state into coordinated, safe, and assistive responses in shared physical environments. 
Because all three stages must run under the power, compute, and reliability constraints of body-worn hardware, lightweight on-device inference is a cross-cutting concern that recurs throughout.


\begin{figure*}[!t]
    \centering
    \includegraphics[width=\linewidth]{figures/part4/section9/wearable_overview.pdf}
    \caption{The wearable AI system as a body-centric loop: on-body signals are modeled into kinematic and biomechanical state (T1), interpreted into activity and intent (T2), and acted on through haptic, exoskeleton, and safety channels (T3), used separately or together. Device and prediction illustrations are schematic.}
    \label{fig:p4_s9_3_overview}
\end{figure*}


\subsection{Pivotal Task-1: Embodied Motion Sensing and Pose Reconstruction}
\label{subsec:p4-sec9-3-pivotal-task-1-embodied-motion-sensing-and-pose-reconstruction}

The first stage of the wearable loop turns raw body-centric signals into a geometric description of how the wearer moves. 
It is a bottom-up reconstruction problem: per-segment inertial and biomechanical measurements must yield joint angles, segment velocities, and ultimately a full-body pose that remains stable over long horizons and in settings where exteroceptive capture fails. 
This task subsumes full-body skeletal pose estimation, fine-grained limb tracking, real-world locomotion and gait analysis, and body-centric localization in GPS-denied spaces; the underlying pose-estimation formulations are developed in Chapter~\ref{chapter:3:sec:6}, so the emphasis here is on reconstruction from on-body sensing. 
Figure~\ref{fig:p4_s9_3_motion} illustrates this reconstruction stack from segment-level signals to full-body pose.

\begin{figure}[!t]
    \centering
    \includegraphics[width=\linewidth]{figures/part4/section9/wearable_motion.pdf}
    \caption{Motion sensing and pose reconstruction from body-worn IMUs: estimators supported by a parametric body model recover root-relative full-body pose, while visual--inertial fusion anchors a world-frame trajectory. Contact or zero-velocity updates apply only when their assumptions hold, and metric scale requires sufficient observability; signals and trajectories are illustrative.}
    \label{fig:p4_s9_3_motion}
\end{figure}

\textbf{From signals to kinematics.} Reconstruction begins with kinematic and dynamic modeling, which provides the physical grounding for interpreting inertial and pressure-based signals. 
Forward and inverse kinematics infer joint angles and segment orientations from inertial measurements, compensating for drift and sensor noise through optimization-based filtering and biomechanical constraints. Classical recursive estimators such as extended and unscented Kalman filters incorporate kinematic priors and temporal smoothness to improve stability in real-world settings~\cite{welch1995kalman, madgwick2011estimation}, while articulated-body dynamics and parametric body models such as SMPL~\cite{loper2015smpl} enforce anatomical plausibility for complex motions such as rapid locomotion or limb coordination. 
Multi-IMU fusion against musculoskeletal models extends these ideas to full-body pose estimation in challenging environments where optical motion-capture systems fail~\cite{huang2018dip, yi2021transpose, yi2022pip}, and recent systems fold in physics-based dynamics and inter-sensor ultra-wideband ranging to tighten accuracy and suppress jitter~\cite{armani2024ultra}.

\textbf{Learning-based reconstruction and sensor fusion.} Where analytic models struggle with sparse or noisy sensing, data-driven approaches learn the mapping from raw inertial streams to pose directly~\cite{huang2018dip, yi2021transpose}, and sensor fusion integrates heterogeneous modalities for drift-compensated estimates. IMU-and-vision fusion resolves ambiguities in inertial tracking by leveraging environmental context from egocentric cameras~\cite{grauman2022ego4d, adt}, while drift-compensation techniques, including zero-velocity updates and constraint-based corrections, stabilize long-horizon wear. 
For body-centric localization, sparse-sensor systems couple inertial pose estimation with simultaneous localization and mapping, so that the wearer's global trajectory is recovered together with body pose rather than left to drift freely~\cite{yi2023egolocate}; visual-inertial odometry that fuses the same inertial signals with first-person imagery offers a complementary route through GPS-denied indoor spaces, assessed on dedicated benchmarks~\cite{Schubert2018TUMVI}.

\textbf{Benchmarks and evaluation.} The kinematic, learning-based, and fusion methods above are measured against the same yardsticks: large-scale motion-capture corpora such as AMASS~\cite{mahmood2019amass}, which provides a canonical surface-and-skeleton reference for full-body pose, alongside dedicated inertial-pose datasets for sparse-IMU reconstruction~\cite{trumble2017totalcapture, huang2018dip}, while body-centric localization is assessed separately on visual-inertial odometry benchmarks~\cite{Schubert2018TUMVI}. Reconstruction quality is reported as mean per-joint position error (MPJPE) and joint-angle error, and localization as absolute trajectory error and accumulated drift (Table~\ref{tab:p4_s9_3_benchmarks}). 
Read against the needs of real deployment, these numbers reveal a clear gap. On curated benchmarks, sparse-IMU systems now recover \emph{root-relative} full-body pose to within roughly five centimeters of optical ground truth~\cite{yi2022pip}, but global translation still drifts over long sequences unless it is anchored by SLAM or ranging~\cite{yi2023egolocate}, and even the root-relative accuracy falls off in the wild, where fast or unseen motions and loose sensor placement push errors past the tolerance that downstream control or clinical gait analysis can absorb. The pose the next stage consumes is thus dependable on average yet brittle at exactly the fast, atypical movements that matter most.


\subsection{Pivotal Task-2: Behavioral Understanding and Intention Modeling}
\label{subsec:p4-sec9-3-pivotal-task-2-behavioral-understanding-and-intention-modeling}

Once motion has been reconstructed, the second stage of the loop interprets it: recognizing what the wearer is doing and anticipating what they will do next. 
Continuous neuromuscular and inertial cues feed activity recognition and segmentation and then short-horizon intention prediction, even in daily-life settings where visual systems face severe occlusion or field-of-view limits. 
Because the recognition and reasoning formulations themselves are treated in Chapter~\ref{sec:reasoning}, the focus here is on how high-frequency wearable streams are modeled to support behavioral inference. 
Figure~\ref{fig:p4_s9_3_behavior} sketches the path from multimodal time-series and body-graph representations to recognized activity and predicted intent.

\begin{figure}[!t]
    \centering
    \includegraphics[width=\linewidth]{figures/part4/section9/wearable_behavior.pdf}
    \caption{Behavioral understanding and intention modeling: IMU and EMG sequences and a kinematic body graph pass through temporal and graph encoders, fused by attention to recognize current activity and forecast short-horizon intent. The EMG pre-activation inset draws anticipatory cues from the observed stream, not from future measurements.}
    \label{fig:p4_s9_3_behavior}
\end{figure}

\textbf{Sequence models for wearable time-series.} The availability of large-scale wearable recordings has driven deep models tailored to high-frequency temporal streams from IMU, EMG, and multimodal sensors. Recurrent networks, temporal convolutional networks, and Transformers are widely applied to gait classification, gesture recognition, and activity segmentation~\cite{ordonez2016deep}: temporal convolutions capture multi-scale dependencies for long-horizon pattern detection~\cite{bai2018tcn}, while attention-based architectures learn complex correlations across sensor channels~\cite{attention}. Because labeled wearable data are scarce and expensive, self-supervised pretraining on unlabeled motion recordings has become the dominant route to transferable representations~\cite{haresamudram2021contrastive}: masked-reconstruction encoders such as LIMU-BERT learn compact IMU features from raw streams~\cite{xu2021limubert}, and multimodal contrastive objectives such as IMU2CLIP align inertial signals with video and language, pointing toward wearable foundation models that generalize across tasks and users~\cite{moon2023imu2clip}.

\textbf{Graph-based body representations.} Because human anatomy is structured as a kinematic tree, graph neural networks offer a natural representation for wearable-derived motion~\cite{kipf2017semisupervised}, encoding relational dependencies between joints and limbs directly in graph space. 
Such formulations fuse IMU placements across body segments, propagate temporal features along skeletal edges, and remain robust to missing or corrupted sensor channels; combining graph convolution with temporal attention jointly models spatial structure and temporal evolution for activity recognition and biomechanical analysis~\cite{yan2018stgcn}.

\textbf{Neuromuscular cues and multimodal fusion for intent.} Beyond recognizing observed motion, EMG-based wearables detect subtle neuromuscular activations that precede visible movement~\cite{atzori2014ninapro}, providing predictive cues for early intention estimation and responsive human-robot collaboration.
Fusing EMG with inertial data sharpens fine-grained gesture and anticipatory-command decoding~\cite{atzori2016deep}, while fusing IMU with egocentric video grounds activity recognition in environmental context~\cite{stein2013combining, grauman2022ego4d}.
A recent milestone shows how far this can scale: a surface-EMG wristband trained across more than a thousand users decodes gestures, clicks, and handwriting from neuromuscular signals alone and, critically, generalizes to new wearers without per-user calibration~\cite{kaifosh2025neuromotor}, addressing the inter-user variability that has long limited EMG interfaces.
Physiological and pressure cues additionally support inference of fatigue, balance, and abnormal gait over extended periods.

\textbf{Benchmarks and evaluation.} Taking the reconstructed motion of the previous stage as input, the sequence, graph, and fusion models above are evaluated on wearable human-activity-recognition corpora~\cite{reiss2012pamap2, chavarriaga2013opportunity}, EMG gesture datasets~\cite{atzori2014ninapro}, and egocentric activity datasets such as Ego4D~\cite{grauman2022ego4d} and EPIC-KITCHENS~\cite{EPIC-KITCHENS}, with recognition reported as accuracy and macro-F1 and intention prediction as anticipation accuracy and prediction lead time (Table~\ref{tab:p4_s9_3_benchmarks}). 
Here the gap to maturity is one of generalization rather than peak score. Balanced locomotion corpora such as PAMAP2 are close to saturated, yet the same architectures score markedly lower on cluttered, null-dominated sets such as Opportunity and lose accuracy again under cross-user and cross-environment evaluation; reliable intention forecasting, meanwhile, still extends only to short, well-rehearsed horizons. This falls well short of the open-world, anticipatory understanding that the safe-collaboration stage below depends on, where the system must read unfamiliar users in unfamiliar settings.


\subsection{Pivotal Task-3: Safe Collaboration and Embodied Interaction}
\label{subsec:p4-sec9-3-pivotal-task-3-safe-collaboration-and-embodied-interaction}

The third stage closes the loop, converting interpreted human state into coordinated, safe, and assistive responses in shared physical environments. 
Contact- and force-level sensing of human--environment interaction supports real-time inference of effort and hazard, which in turn drives feedback for collaboration and shared autonomy. 
The collaboration and shared-control reasoning that this task builds on is developed in Chapter~\ref{sec:reasoning} and, for the manipulation setting specifically, in Section~\ref{sec:p4_s9_1}; here the emphasis is on the wearable feedback and safety layer. 
Figure~\ref{fig:p4_s9_3_safety} depicts this closing stage, from contact sensing to coordinated safety response.

\begin{figure}[!t]
    \centering
    \includegraphics[width=\linewidth]{figures/part4/section9/wearable_safety.pdf}
    \caption{Safe human--robot collaboration with wearables: fused contact, force, and proximity signals, together with intent forecasts (T2) and robot motion context, drive motion prediction and risk evaluation, which select glove, exoskeleton, or cobot responses. Prediction envelopes and risk thresholds are schematic, not measured safety guarantees.}
    \label{fig:p4_s9_3_safety}
\end{figure}

\textbf{Sensing interaction and effort.} Haptic and tactile wearables capture contact forces, surface interactions, and manipulative behaviors, supplying the signals needed to model human--environment interaction and object handling~\cite{HOI4D, luo2021learning}. Exoskeletons equipped with torque and force sensors use these signals to infer user effort and detect hazardous conditions, adjusting assistive strategies accordingly~\cite{young2017state}; probabilistic fusion that integrates biomechanical priors and uncertainty estimates is central to robustness in such safety-critical settings.

\textbf{Coordination, safety, and feedback.} In industrial and shared-autonomy scenarios, on-body inertial and proximity sensors contribute to safety monitoring by detecting collision risks, unsafe postures, and sudden motion anomalies, while haptic channels return spatial alerts and force feedback that close the sensorimotor loop~\cite{pacchierotti2017wearable}. 
Combined with the intention forecasts of Pivotal Task-2, these signals form the basis for predictive safety mechanisms and shared-control frameworks that anticipate human motion rather than merely reacting to it.

\textbf{On-device constraints.} Because feedback must be timely, the action stage is where the wearable system's power, memory, and compute limits bite hardest. Model compression (pruning, quantization, and distillation)~\cite{hinton2015distilling} and edge-optimized architectures~\cite{howard2017mobilenets} enable real-time inference on wearable microcontrollers or smartphone-class processors, and on-device adaptation mitigates inter-user variability without relying on remote servers, keeping safety-critical computation local and low-latency.

\textbf{Benchmarks and evaluation.} Closing the loop, the interaction-sensing, safety, and on-device methods above are assessed on human-object-interaction datasets such as HOI4D~\cite{HOI4D} and, for assistive control, on exoskeleton and teleoperation testbeds, with metrics spanning contact- and force-estimation error, collision-avoidance and safety-violation rates, and intention-anticipation lead time (Table~\ref{tab:p4_s9_3_benchmarks}). 
Evaluation here is the least mature of the three stages: it remains fragmented across hardware-specific testbeds with few shared protocols, so reported gains are hard to compare and seldom stress real safety guarantees. 
Current systems perform convincingly in instrumented demonstrations, but closed-loop, certified safety under the full variability of human partners, the bar that genuine shared autonomy must clear, is still an open target rather than a measured result. The robustness, generalization, and safety gaps that recur at each stage of the loop are what the challenges below take up directly.


\begin{table}[!th]
\centering
\caption{Representative benchmarks for wearable spatial-intelligence tasks, grouped by pivotal task.}
\label{tab:p4_s9_3_benchmarks}
\fontsize{8}{10}\selectfont
\setlength{\tabcolsep}{1.0mm}
\begin{tabular}{L{0.16\textwidth} L{0.20\textwidth} L{0.34\textwidth} L{0.21\textwidth}}
\toprule
\textbf{Pivotal Task} & \textbf{Benchmark / Dataset} & \textbf{Representative method / use} & \textbf{Headline metric} \\
\midrule
T1: Motion \& Pose & AMASS~\cite{mahmood2019amass} & Full-body pose reference (SMPL surface) & MPJPE \\
T1: Motion \& Pose & TotalCapture~\cite{trumble2017totalcapture} & Video+inertial full-body pose & MPJPE \\
T1: Motion \& Pose & DIP-IMU~\cite{huang2018dip} & Sparse-IMU pose reconstruction (DIP, TransPose, PIP) & Joint-angle / position error \\
T1: Localization & TUM-VI~\cite{Schubert2018TUMVI} & Visual-inertial odometry (body-centric) & ATE / drift \\
T2: Activity & PAMAP2~\cite{reiss2012pamap2} & Wearable activity recognition (deep recurrent / TCN) & Accuracy / macro-F1 \\
T2: Activity & Opportunity~\cite{chavarriaga2013opportunity} & On-body activity recognition (DeepConvLSTM) & Weighted / macro-F1 \\
T2: Activity & Ego4D~\cite{grauman2022ego4d} & Egocentric activity recognition & Accuracy / F1 \\
T2: Gesture & Ninapro~\cite{atzori2014ninapro} & EMG gesture / movement classification & Accuracy \\
T2: Gesture & emg2qwerty~\cite{sivakumar2024emg2qwerty} & Large-scale surface-EMG typing & Character error rate \\
T3: Interaction & HOI4D~\cite{HOI4D} & Human-object interaction sensing & Contact / interaction acc. \\
\bottomrule
\end{tabular}
\end{table}


\subsection{Established Tools and Systems}
\label{subsec:p4-sec9-3-established-tools-and-systems}
Commercial wearables have converged on the head and the wrist, and what separates current products is less their sensors than where perception runs and whether anything is shown back to the wearer. Full augmented-reality headsets are covered with the XR platforms in Section~\ref{sec:p4_s10_1}; here the focus is on everyday-wearable glasses, neural input bands, and the research platforms and chips behind them.

\textbf{AI glasses without a display.} By volume, the dominant wearable AI device is a pair of ordinary-looking glasses carrying a camera, microphones, open-ear speakers, and an IMU, with a multimodal assistant running on the paired phone or in the cloud. Meta's line with EssilorLuxottica\footnote{\url{https://www.meta.com/ai-glasses/}} defines the category: the second-generation Ray-Ban Meta (2025) records 3K video from a 12-megapixel camera and runs up to eight hours per charge, and the Oakley Meta Vanguard adapts the design to sport with a 122-degree camera. Display-less models made up 96\% of AI-glasses shipments in the first half of 2026, almost all of them Meta's\footnote{\href{https://counterpointresearch.com/en/insights/global-ai-glasses-shipments-surge-263-percent-yoy-with-display-less-ai-glasses}{Counterpoint Research, global AI glasses shipments, H1 2026}.}, while Xiaomi's AI glasses (2025) and Alibaba's Qwen-based Quark AI glasses\footnote{\url{https://www.quark.cn/}} (November 2025) lead a fast-growing Chinese segment. Google and Samsung previewed Gemini-based Android XR glasses\footnote{\url{https://www.android.com/xr/}} with Gentle Monster and Warby Parker in May 2026, with the first display-less models due in the fall. Spatially, these devices answer questions about what the wearer is looking at one camera frame at a time; most keep no persistent map of the surroundings, so the continuous body and scene tracking of Pivotal Tasks~1 and~2 is still largely absent from shipping products.

\textbf{Display glasses and neural input.} Adding a display makes the glasses an interface, and input then becomes the hard problem. Meta Ray-Ban Display\footnote{\url{https://www.meta.com/ai-glasses/meta-ray-ban-display/}} (September 2025) puts a small monocular display, with a 20-degree field of view, in the right lens and ships with the Meta Neural Band, a surface-EMG wristband that decodes finger gestures and handwriting; it is the commercial form of the generic neuromotor interface described in Pivotal Task-2~\cite{kaifosh2025neuromotor}. In waveguide display glasses, Rokid\footnote{\url{https://global.rokid.com/}} led shipments in the first half of 2026, ahead of Meta and Even Realities\footnote{\url{https://www.evenrealities.com/}}, whose glasses favor a minimal text display over cameras.

\textbf{Research platforms and on-glass compute.} Much of the egocentric data behind the methods of this section was recorded on research glasses rather than consumer products. Meta's Project Aria~\cite{engel2023projectaria}\footnote{\url{https://www.projectaria.com/}} provides glasses with RGB, tracking, and eye-tracking cameras, IMUs, and spatial microphones, together with machine-perception services for trajectories, eye gaze, and hand tracking; its second generation (2025) adds a heart-rate sensor and a contact microphone in the nosepad, alongside on-device perception. The platform underlies Aria Digital Twin~\cite{adt}; Ego-Exo4D~\cite{grauman2024egoexo4d}, which pairs egocentric Aria recordings of skilled activities with synchronized exocentric cameras over 1{,}422 hours; and Nymeria~\cite{ma2024nymeria}, which adds full-body ground-truth motion from a motion-capture suit to 300 hours of everyday activity, the kind of supervision the pose reconstruction of Pivotal Task-1 needs. Ego4D~\cite{grauman2022ego4d} gathered its 3{,}670 hours from several head-mounted cameras, and remains the main benchmark for egocentric activity understanding. On the compute side, Qualcomm's Snapdragon AR1+ Gen~1\footnote{\url{https://www.qualcomm.com/products/mobile/snapdragon/xr-vr-ar}} (2025) runs a one-billion-parameter language model on the glasses without a phone, a first step toward the on-device inference that the latency and power constraints of Pivotal Task-3 call for.

\subsection{Challenges}
\label{subsec:p4-sec9-3-challenges}

Wearable sensing is coupled to the dynamics and variability of the human body in a way that external sensing is not, and that coupling is the source of its hardest problems. The four below track the loop of the System Overview: signal quality at the sense stage, robustness under partial observation, fusion across heterogeneous modalities, and generalization across users.

\paragraph{Signal Instability and Drift.}
On-body signals are inherently unstable. IMU drift, EMG noise, and calibration shifts arise continuously from motion artifacts, sensor displacement, and soft-tissue deformation, and they strike precisely the quantities that pose estimation, trajectory tracking, and dynamics reconstruction integrate over time, so small per-step errors compound into large ones across a long recording.

\paragraph{Incomplete and Intermittent Sensing.}
Sensing is rarely continuous in practice. Sweat, motion, and wireless dropouts cause brief losses of contact and stretches of degraded signal, leaving the system to reconstruct reliable motion from partial and sometimes corrupted segment-level data, a demand that clean-benchmark evaluation rarely exercises.

\paragraph{Heterogeneous Multimodal Fusion.}
The modalities that make wearables powerful are also mismatched. IMU, EMG, pressure, and egocentric-visual streams differ in sampling rate, latency, noise, and coordinate frame, and fusing them under real-world variability, without letting one noisy channel corrupt the whole estimate, remains an open integration problem rather than a solved preprocessing step.

\paragraph{Inter-User Variability.}
Finally, the body itself is the confound. Anatomy, muscle-activation patterns, and movement style vary enough between people that a model tuned on a small or homogeneous population transfers poorly, and the effect is sharpest exactly where wearables matter most, in motion reconstruction, intent prediction, and gait analysis, where individual biomechanics dominate the sensor signature. Personalization, domain adaptation, and larger, more diverse training corpora are the levers for closing this gap, and how far they succeed will decide whether wearable spatial intelligence moves from convincing demonstrations to dependable deployment.

\newpage
\section{Aerial Agents}
\label{sec:p4_s9_4}


% AeroVerse: UAV-Agent Benchmark Suite for Simulating, Pre-training, Finetuning, and Evaluating Aerospace Embodied World Models
% Airvista-ii: An agentic system for embodied uavs toward dynamic scene semantic understanding
% Towards Autonomous UAV Visual Object Search in City Space: Benchmark and Agentic Methodology
% Next-Generation LLM for UAV: From Natural Language to Autonomous Flight

\subsection{Introduction}
\label{subsec:p4-sec9-4-introduction}


Aerial agents (satellites, high-altitude platform stations (HAPS), and unmanned aerial vehicles (UAVs)) acquire rich, dynamic, large-scale observations from the sky and near space and, increasingly, perceive, decide, and act on them autonomously through onboard processing and adaptive control (Table~\ref{tab:air-agents}).
The aerial domain is a demanding testbed for the spatial capabilities surveyed in Part~\ref{part:3}: it stresses perception across extreme scale and viewpoint changes (Chapter~\ref{chapter:3:sec:6}) and decision-making under partial observability and tight latency budgets (Chapter~\ref{sec:reasoning}), across platform tiers that differ in altitude, resolution, endurance, and latency by orders of magnitude.

Satellites supply the global tier. Navigation constellations, namely GPS~\cite{enge1994global,leick2015gps} within the broader Global Navigation Satellite System (GNSS), Russia's GLONASS~\cite{ivanov1992glonass}, the European Union's Galileo, and China's BeiDou~\cite{yang2019introduction,yang2020basic}, provide the positioning, navigation, and timing on which vehicles~\cite{li2022real}, smartphones~\cite{zangenehnejad2021gnss}, network synchronization~\cite{hasan2018gnss}, and every lower aerial tier rely.
Earth-observation satellites operate continuously over years, which suits long-term, wide-area monitoring that neither ground sensors nor short-endurance UAVs can provide: seasonal vegetation dynamics~\cite{malingreau1986global}, wildlife migration~\cite{dobrynin2017integration}, river-basin ecosystem health~\cite{ferreira2013estimating}, deforestation~\cite{tucker2000strategies,finer2018combating}, desertification~\cite{del1998status,collado2002satellite}, biodiversity loss~\cite{turner2015free}, and climate indicators such as land surface temperature, ocean color, and atmospheric composition~\cite{yang2013role,kaufman2002satellite}, together with ecological conservation~\cite{pettorelli2014satellite,boyle2014high}, multispectral and hyperspectral resource discovery~\cite{pettorelli2019satellite}, precision agriculture~\cite{guo2018multi} and crop-yield estimation~\cite{maas1988using,ferencz2004crop}, and the detection of algal blooms~\cite{fang2022global} and pollutant leakage~\cite{pandey2019satellite}.

\begin{table}[!ht]
\centering
\caption{Comparison of satellite, HAPS, and UAV platforms.} 
\label{tab:air-agents}
\small
\setlength{\tabcolsep}{2.5mm}
\begin{tabular}{@{}llll@{}}
\toprule
\textbf{Parameter} & \textbf{Satellite} & \textbf{HAPS} & \textbf{UAV} \\
\midrule
Altitude & 500--36,000 km & 20--50 km & < 10 km \\
Spatial Resolution & m to sub-m & dm to cm & cm-level \\
Temporal Coverage & Hours--days revisit & Persistent (regional) & Real-time (local) \\
Coverage Area & Global/regional & Regional & Localized \\
Endurance & Years & Months & Hours \\
Maneuverability & Minimal & Limited & High \\
Data Latency & High & Medium & Low \\
Payload Capacity & High & Moderate & Low \\
Onboard Processing & Limited & Moderate & High \\
Atmospheric Impact & Minimal & Low & High \\
Deployment Speed & Slow & Medium & Fast \\
Cost Structure & High (low per area) & Medium & Low (high per area) \\
%AI Suitability & Large-scale datasets & Regional monitoring & Real-time inference \\
\bottomrule
\end{tabular}
\end{table}


Differently, UAVs, commonly known as drones, have gained widespread popularity due to their agility, low operational cost, and capability for autonomous operation. Equipped with onboard sensors and processing units, UAVs can perceive their surroundings, plan flight paths, and adapt to dynamic conditions in real time. This unique perspective allows UAVs to capture high-resolution imagery and video from flexible vantage points, enabling frequent revisits and reduced occlusions in complex environments. Compared to satellites, UAVs excel in capturing fine-grained, localized data, and can be rapidly deployed to areas of interest, including hazardous or otherwise inaccessible regions. They now serve aerial photography~\cite{zhao2023ms,shao2022efficient}, cinematography~\cite{pitas2022autonomous,dang2022path}, environmental mapping~\cite{zheng2020university,wang2021each,lin2022joint}, precision agriculture~\cite{goodrich2023placement}, geolocalization~\cite{liu2019lending,shi2019spatial,shi2020optimal,shi2020looking,shi2022beyond,shi2022accurate,shi2022cvlnet,shi2023boosting,song2023learning}, and last-mile delivery~\cite{sorbelli2023wind,dissanayaka2023review}.

Filling the altitude gap between satellite systems and UAVs, emerging platforms such as stratospheric balloons, solar-powered high-altitude platform stations, and small satellite constellations (\eg, CubeSats~\cite{lemmer2017propulsion}) are increasingly being integrated into aerial observation networks.
These high-altitude platform stations (HAPS) offer complementary advantages: HAPS~\cite{gonzalo2018capabilities} provide persistent regional coverage with higher spatial resolution than satellites, while orbital systems~\cite{cunningham2004national} enable global, long-term monitoring at varying temporal and spectral scales.
Together, they form a multi-tiered aerial ecosystem capable of capturing heterogeneous multimedia data, from visible and thermal imagery to hyperspectral and LiDAR scans, across spatiotemporal scales unattainable by any single platform alone.
For instance, solar-powered HAPS such as Airbus Zephyr~\cite{kothe2019flight} or HAWK30 by AeroVironment/HAPSMobile~\cite{fladeland2019demonstrating} can operate autonomously in the stratosphere (about 20 km altitude) for weeks, offering persistent observation and communication relay services with sub-meter resolution, which is superior to most Earth observation satellites and at lower cost than maintaining UAV fleets.

A growing trend is the coordination of heterogeneous aerial agents into cooperative networks, where satellites offer global context, HAPS serve as communication relays or edge computing nodes, and UAVs execute localized missions, enabling real-time, closed-loop intelligence. The utility of these systems is further amplified by advances in artificial intelligence, including machine learning, computer vision, big data analytics, and transfer learning~\cite{wang2024multiple,wang2022learning,zheng2022adaptive}. These technologies empower intelligent processing of massive aerial datasets, enabling automated image and video analysis~\cite{liu2023robust,li2021uav,9854892,zheng2019unsupervised,wang2023context,ju2025video2bev}, real-time object detection and tracking~\cite{zhang2019eye,wen2021detection}, predictive modeling for urban dynamics~\cite{fang2022cognitive,shi2022geometry}, natural language-based control interfaces~\cite{chu2024towards}, and city-scale digital twin simulations~\cite{zheng2022citysim}. Data privacy, airspace integration, energy efficiency, and operational robustness remain open, and are taken up in the Challenges below.


\subsection{System Overview}
\label{subsec:p4_s9_4_overview}
An aerial agent realizes a sensing-to-action loop that spans the multi-tier ecosystem of Table~\ref{tab:air-agents}: satellites supply global context and positioning, HAPS act as persistent regional relays and edge-computing nodes, and UAVs execute localized, low-latency missions at the point of interest. Across this stack, spatial intelligence is exercised through three complementary stages, which constitute the pivotal tasks of this section. Cross-view perception anchors aerial and ground observations to a common geographic frame, so that what a drone sees can be matched to a map or a street-level query (Pivotal Task-1); remote-sensing detection and environmental understanding turn that imagery into structured scene knowledge about objects, layouts, and human activity (Pivotal Task-2); and autonomous navigation and decision making convert perception into flight plans and control actions in unseen environments (Pivotal Task-3). Figure~\ref{fig:p4_s9_4_overview} sketches this pipeline, with a feedback path from action back onto sensing that ties the three stages into one closed loop rather than a feed-forward chain.

\begin{figure}[!ht]
    \centering
    \includegraphics[width=0.98\linewidth]{figures/part4/section9/aerial_overview.pdf}
    \caption{Spatial intelligence for aerial agents as a closed loop from multi-tier aerial sensing (satellite, HAPS, UAV, GNSS, onboard sensors), through three complementary core capabilities, to flight plans with 6-DoF pose targets and downstream flight control. The controller executes the plan; subsequent observations close the sensing loop. The core comprises cross-view geolocalization, detection and environmental understanding, and navigation and decision making. Land-use, water, and infrastructure layers are parallel examples. Platform and scene imagery is illustrative.}
    \label{fig:p4_s9_4_overview}
\end{figure}

\subsection{Pivotal Task-1: Ground-Air Collaborative Perception and Visual Geolocalization}
\label{subsec:p4_s9_4_task1}

A foundational requirement of any aerial system is to place its observations on the map: the integration of ground-level and aerial views for cross-modal understanding enables visual geolocalization, determining the geographic location of a query image by matching it against a reference database. This task anchors the perception stage of the system and is particularly challenging due to the large modality gap between ground street-view images and overhead (aerial or satellite) imagery, an instance of the cross-view correspondence problem treated more generally in Chapter~\ref{chapter:3:sec:6}.
For instance, the \textit{University-1652} dataset~\cite{zheng2020university} has become a benchmark for this task, offering ground, drone and satellite image triplets of campus buildings, thus facilitating research on cross-view matching and localization. Such capabilities are critical for urban search-and-rescue, augmented reality, and infrastructure monitoring, where precise alignment between ground observations and aerial maps is essential.
Progress has come from bridging the modality gap, encoding spatial structure into a viewpoint-invariant representation, and exploiting multiple views per location (Figure~\ref{fig:p4_s9_4_geoloc}).



\textbf{Bridging the modality gap.} Most extant research frames geolocalization as a cross-view image retrieval problem, where the core challenge is to learn a viewpoint-invariant representation that effectively bridges the disparity between images captured from different perspectives. Recent advances in deep learning have led to the widespread application of Convolutional Neural Networks (CNNs) and Vision Transformers (ViTs) for extracting robust visual features. One prominent line of inquiry employs metric learning to construct a shared latent space for images acquired from diverse platforms. Workman~\etal~\cite{workman2015location} show that a CNN pre-trained on the Places dataset~\cite{zhou2014learning} can, by itself, yield highly discriminative features without explicit fine-tuning for the task. The integration of contrastive loss, which minimizes the distance between positive pairs in the feature space, further enhances geolocalization performance~\cite{workman2015wide,lin2015learning}. Liu~\etal~\cite{liu2019stochastic} introduce the Stochastic Attraction and Repulsion Embedding (SARE) loss, which minimizes the Kullback-Leibler divergence between the learned and actual data distributions, and more recently Deuser~\etal~\cite{deuser2023sample4geo} show that hard negative sampling strategies also significantly influence the final discriminative power.

\textbf{Encoding spatial structure.} Another significant body of work addresses the problem of spatial misalignment in ground-to-aerial matching, injecting geometric priors into the representation. Vo~\etal evaluate various network architectures and propose an orientation regression loss to train an orientation-aware network~\cite{vo2016localizing}. Zhai~\etal~\cite{zhai2017predicting} utilize semantic segmentation maps to facilitate semantic alignment, while Hu~\etal~\cite{hu2018cvm} incorporate a NetVLAD layer to extract more discriminative features. Subsequently, Liu~\etal~\cite{liu2019lending} propose a Siamese Network to explicitly integrate spatial cues, such as orientation maps, into the training process, and Shi~\etal~\cite{shi2020optimal} introduce a spatial-aware layer to improve localization performance. Given that each location has multiple training images from different viewpoints, Zheng~\etal~\cite{zheng2020university} train a multi-branch CNN as a base model. Building upon this multi-branch architecture, Wang~\etal~\cite{wang2021each} further partition the feature map into multiple segments to improve matching discriminability, while Lin~\etal~\cite{lin2022joint} incorporate keypoint detection into the learning process. More recently, researchers resort to ViT structures: L2LTR~\cite{yang2021cross} adopts a hybrid CNN-ViT architecture, leveraging CNNs to extract high-level semantic features and ViTs to encode them into global representations; GeoDTR~\cite{zhang2023cross} further enhances geometric reasoning by introducing a ViT-based layout extractor that generates geometry-aware attention maps; and TransGeo~\cite{zhu2022transgeo} advances the paradigm with a two-branch ViT model and a two-stage coarse-to-fine learning strategy.

\begin{figure}[!ht]
    \centering
    \includegraphics[width=0.98\linewidth]{figures/part4/section9/aerial_geoloc.pdf}
    \caption{Cross-view visual geolocalization as viewpoint-invariant metric learning: the illustrated shared encoder embeds ground, drone, and satellite views so that same-place (positive) pairs cluster and different-place negatives separate. Architecture variants may instead use view-specific branches. Arrows within the feature-space sketch indicate attraction and repulsion during representation learning, not geographic motion. Synthetic views illustrate a common site rather than a calibrated image triplet.}
    \label{fig:p4_s9_4_geoloc}
\end{figure}

\begin{table}[t]
\centering
\setlength{\tabcolsep}{3pt}
\caption{Comparison of visual geolocalization datasets. R@K: recall at top $K$; AP: average precision; PR: precision-recall.}
\label{table:geo-dataset}
\small
\resizebox{\linewidth}{!}{
\begin{tabular}{lccccccc}
\toprule
Dataset & Univ.-1652~\cite{zheng2020university} & SUES-200~\cite{zhu2023sues} & CVUSA~\cite{zhai2017predicting} & CVACT~\cite{liu2019lending} & Lin~\etal~\cite{lin2015learning} & Tian~\etal~\cite{tian2017cross} & Vo~\etal~\cite{vo2016localizing} \\
\midrule
\#Train imgs. & 701$\times$71.6 & 120 loc. & 35.5k$\times$2 & 35.5k$\times$2 & 37.5k$\times$2 & 15.7k$\times$2 & 900k$\times$2 \\
Platform & Drone, Ground, Sat. & Drone, Sat. & Ground, Sat. & Ground, Sat. & Ground, 45$^{\circ}$ Aerial & Ground, 45$^{\circ}$ Aerial & Ground, Sat. \\
Imgs./loc. & 54+16.6+1 & drone + 1 & 1+1 & 1+1 & 1+1 & 1+1 & 1+1 \\
Target & Building & Diverse & User & User & Building & Building & User \\
Eval. & R@K, AP & R@K, AP & R@K & R@K & PR, AP & PR, AP & R@K \\
\bottomrule
\end{tabular}}
\end{table}

\textbf{From image pairs to multi-view benchmarks.} The maturation of these methods is closely tied to the datasets that score them (see Table~\ref{table:geo-dataset}). The majority of prior geolocalization datasets are based on image pairs, designed for matching images from two distinct platforms such as terrestrial cameras and satellites. One of the earliest contributions~\cite{lin2015learning} leverages public sources to create 78,000 image pairs spanning a 45-degree bird's-eye perspective and the ground. In a similar vein, Tian~\etal~\cite{tian2017cross} compile image pairs for urban localization, positing that buildings are crucial for this task and thus incorporating building detection into their pipeline. The CVUSA~\cite{zhai2017predicting} and CVACT~\cite{liu2019lending} datasets instead address matching panoramic ground-view images with satellite-view images, which is particularly useful for user localization when GPS data is unavailable. A key distinction is one of objective: the latter two focus on localizing the user who captures the photograph, while the former focus on localizing a specific target within the image, so that multiple views of the target are more beneficial for the target-localization datasets (Lin~\etal, Tian~\etal) because they encourage the model to learn the target's structure. However, most existing datasets typically provide only two views of a given location. The University-1652 dataset~\cite{zheng2020university} addresses this by including a greater number of views per target, thereby fostering more robust viewpoint-invariant features; SUES-200~\cite{zhu2023sues} follows a similar multi-view strategy with drone-captured real-world data. Wang~\etal~\cite{wang2024multiple} further extend University-1652 to multi-weather evaluation (rain, fog, and other complex conditions), while Chu~\etal~\cite{chu2024towards} introduce linguistic annotations and a language-based query setting for a more intuitive user experience.

\textbf{Metrics.} These benchmarks are scored by retrieval metrics that reflect their gallery structure. Several datasets~\cite{liu2019lending,zhai2017predicting,vo2016localizing} adopt Recall@K, defined as 1 if a correct match appears within the top K retrieved images; Recall@K is highly sensitive to the rank of the first correct match and is best suited for test sets containing only a single ground-truth match in the gallery. In contrast, datasets such as University-1652~\cite{zheng2020university} and SUES-200~\cite{zhu2023sues} include multiple ground-truth matches from various viewpoints, where Recall@K fails to comprehensively reflect retrieval performance; they therefore also employ Average Precision (AP), as used by Lin~\etal~\cite{lin2015learning} and Tian~\etal~\cite{tian2017cross}. AP, defined as the area under the Precision-Recall (PR) curve, considers all ground-truth images in the gallery, and recent evaluations usually report both Recall@K and the mean Average Precision (mAP) across all queries.


\subsection{Pivotal Task-2: Remote Sensing Detection and Environmental Understanding}
\label{subsec:p4_s9_4_task2}

Once observations are geolocalized, the next stage turns aerial imagery into structured scene knowledge. Aerial platforms are integral to large-scale environmental monitoring and public safety, facilitating fine-grained object detection in complex scenes, a specialization of the detection and segmentation machinery of Chapter~\ref{chapter:3:sec:6} to the aerial viewpoint. Prominent datasets such as \textit{DOTA}~\cite{xia2018dota}, \textit{DOTA-v2}~\cite{xia2018dota,ding2021object}, and \textit{VisDrone}~\cite{du2019visdrone} provide comprehensive annotations for detecting various objects, including vehicles and pedestrians, under a wide range of weather, illumination, and scale conditions. These benchmarks have spurred the development of robust detection models specifically engineered to address challenges inherent to aerial perspectives, such as extreme aspect ratios, dense object clustering, and diminutive object sizes. A significant application is found in UAV-based wildlife monitoring, where systems are deployed for animal census in conservation areas~\cite{vermeulen2013unmanned}, highlighting the utility of aerial technologies in promoting ecological sustainability. Figure~\ref{fig:p4_s9_4_detection} relates the aerial-specific difficulties of orientation, scale, and density to the detectors built for them and to the human-centric understanding that aerial video enables.

\begin{figure}
    \centering
    \includegraphics[width=0.98\linewidth]{figures/part4/section9/aerial_detection.pdf}
    \caption{Remote-sensing detection from the aerial viewpoint: arbitrary orientation, extreme scale variation, and dense small objects (top) motivate (1) oriented-object detection and (2) multi-scale feature learning, which provide complementary representations for human-centric understanding (bottom). Person re-identification matches identities across views; action examples are alternative classes, not a prescribed sequence; tracking maintains identities across frames. Synthetic scenes illustrate these tasks.}
    \label{fig:p4_s9_4_detection}
\end{figure}

Beyond object detection, as shown in Table~\ref{tab:uav_reid_datasets}, with the proliferation of UAVs in smart-city and other human-centric applications, human behavior understanding~\cite{li2021uav} has emerged as a critical and formidable research challenge, including person re-identification and action recognition. For instance, the UAV person re-identification (re-ID) task~\cite{albaluchi2025uav} is designed to address cross-drone identity matching amid substantial viewpoint and illumination variations, thereby motivating the creation of view-invariant feature representations. In conjunction with resources like \textit{VisDrone}, which contains multi-target tracking sequences from mobile UAVs, these datasets facilitate the design of intelligent video analytics systems capable of performing long-term human behavior analysis and anomaly detection in dynamic environments.


\begin{table}[t]
\centering
\caption{Comparison of UAV-based human and object datasets.}
\label{tab:uav_reid_datasets}
\fontsize{8}{10}\selectfont
\setlength{\tabcolsep}{1.0mm}
\begin{tabular}{L{0.15\textwidth} L{0.26\textwidth} L{0.25\textwidth} L{0.25\textwidth}}
    \toprule
    \textbf{Feature} & \textbf{UAV-Human}~\cite{li2021uav} & \textbf{DOTA-v2}~\cite{xia2018dota,ding2021object} & \textbf{VisDrone}~\cite{du2019visdrone} \\
    \midrule

    \textbf{Primary Focus} & Human Behavior Understanding & Object Detection & Object Detection / Tracking / Crowd Counting \\

    \textbf{Data Type} &  Videos \& Images & Images & Videos \& Images \\

    \textbf{\#Images} & 22,476 frames for pose estimation, 41,290 frames for re-ID, 22,263 frames for attribute recognition & 11,268 & 10,209 \\

    \textbf{\#Videos} & 67,428  & N/A & 288 \\

    \textbf{\#Categories} & 155 action categories & 18 object categories & 10 object categories \\

    \textbf{\#Instances} & 119 subjects for action recognition, 1,144 identities for re-ID & 1,793,658 & > 2.6 million annotated bounding boxes \\

    \textbf{Sensor Modalities} & RGB, Depth, Infrared (IR), Fisheye, Night-Vision & RGB (from Google Earth, aerial) \& Grayscale (from satellites) & RGB \\

    \textbf{Image Resolution} & 1920$\times$1080 (RGB), 640$\times$576 (Depth,IR), 640$\times$480(Fisheye, Night-Vision), & 800$\times$800 to 20,000$\times$20,000 pixels & 2000$\times$1500 (image),  3840$\times$2160 (video) \\

    \textbf{Annotation Format} & Bounding boxes, 17 keypoints for pose & Oriented Bounding Boxes (Quadrilaterals) & Bounding boxes with attributes (\eg, occlusion) \\

    \textbf{Scene Diversity} & Urban and rural, daytime and nighttime & Wide area, diverse landscapes (ports, airports, cities, etc.) & 14 different cities, urban and rural, various weather and lighting \\

    \textbf{Platform} & Drone (DJI Matrice 100) & Satellite, Aerial Platforms & Drone \\

    %\textbf{Unique Features} & Rich multi-modal data for fine-grained human-centric analysis. & Very high-resolution images with oriented bounding boxes for rotated objects. & A comprehensive benchmark with multiple tailored tasks for drone vision. \\
    \bottomrule
  \end{tabular}
\end{table}


Remote sensing object detection is primarily challenged by the prevalence of arbitrarily oriented objects and substantial scale variations~\cite{xia2018dota, ding2021object, sun2022fair1m, liu2016ship, cheng2022anchor, yao2023automated}, and existing research can be broadly organized along two corresponding fronts.
\textbf{Detection of arbitrarily oriented objects.} A significant line of work focuses on developing specialized oriented object detection (OOD) methods. Early approaches often augment horizontal bounding boxes (HBBs) with an angle parameter, employing techniques such as feature refinement in the detector neck~\cite{yang2019scrdet,yang2021r3det}, extraction of rotated regions of interest (RoIs)~\cite{ding2019learning,xie2021oriented}, or designing specific OOD heads~\cite{hou2022shape,han2021redet,ming2021dynamic}. However, these methods frequently suffer from boundary discontinuity and representation ambiguity due to their reliance on angle parameters. To mitigate these issues, subsequent research has explored more flexible \textit{object representations}: some methods describe oriented objects using points~\cite{fu2020point, li2022oriented} or gliding vertexes~\cite{xu2020gliding}, while others model them with Gaussian distributions~\cite{yang2021rethinking, yang2021learning, cheng2022dual} and design corresponding loss functions~\cite{qian2021learning}.
\textbf{Handling substantial scale variations.} Another research direction designs specialized backbones and feature extraction strategies to cope with large scale variations in remote sensing images (RSIs). A common strategy is multi-scale feature learning, achieved through data augmentation~\cite{zhao2019multi, shamsolmoali2021rotation, chen2020stitcher}, multi-scale feature integration~\cite{zhang2019hierarchical,zheng2020hynet,9382268,liu2016ssd}, enhanced feature pyramid networks (FPNs)~\cite{lin2017feature,hou2022refined,zhang2021laplacian,guo2020rotational}, or multi-scale anchor generation~\cite{guo2018geospatial, hou2021self, qiu2019a2rmnet}. Other works focus on ensuring equivalent receptive fields for objects of different orientations~\cite{han2021redet,pu2023adaptive}. More recently, some methods~\cite{Li_2023_ICCV} have employed large kernel convolutions~\cite{liu2022convnet,ding2022scaling,liu2022more} to enlarge the spatial receptive field for large objects, though this can introduce background noise for smaller ones, while others harness multi-scale convolution kernels and attention to address scale variance~\cite{guo2022segnext,dai2023multi,zhang2022multiscale,cai2024poly}.

These method fronts are evaluated on a complementary set of benchmarks and metrics. The VisDrone dataset~\cite{du2019visdrone} provides a comprehensive benchmark for drone-captured imagery, featuring urban and rural scenes with challenges like varying altitudes and occlusions; DOTA-v2~\cite{xia2018dota,ding2021object} focuses on oriented object detection in high-resolution satellite images, emphasizing rotated bounding boxes for objects like planes and ships; and UAV-Human~\cite{li2021uav} targets human behavior understanding from UAVs, with multi-modal sequences for pose and action analysis. Complementary datasets, such as UAVDT~\cite{du2018uavdt} for traffic surveillance and iSAID~\cite{waqas2019isaid} for instance segmentation, broaden the scope of aerial detection tasks and collectively enable cross-dataset fine-tuning to mitigate domain gaps. Evaluation extends standard metrics to address orientation, scale, and density: the primary metric, mean Average Precision (mAP), is computed as the area under the precision-recall curve, typically at an Intersection over Union (IoU) threshold of 0.5 (mAP@0.5) or across 0.5:0.95 for stricter assessment~\cite{lin2014microsoft}, and for oriented detection, as in DOTA-v2, rotated IoU variants ensure alignment with object headings~\cite{xia2018dota}. Per-class and per-size Average Precision exposes imbalances, and small objects ($<32\times32$ pixels) remain the main source of error on drone imagery~\cite{du2019visdrone}. Additional metrics include precision, recall, and F1-score to assess false positives (\eg, background clutter) and misses (\eg, occlusions), while Frames Per Second (FPS) is critical for real-time UAV deployment, targeting 30+ FPS~\cite{yang2023laiyolo}. Video-specific metrics like Multiple Object Tracking Accuracy (MOTA) are used for dynamic scenes in datasets like UAV-Human~\cite{li2021uav}. Modern real-time detectors of the YOLO family~\cite{wang2024yolov9} are increasingly adapted to the aerial setting, though aggregated aerial mAP still trails their COCO-level performance owing to small-object and orientation hardness, keeping open gaps such as multi-modal fusion and zero-shot adaptation active.


\begin{figure}[!ht]
    \centering
    \includegraphics[width=0.98\linewidth]{figures/part4/section9/aerial_navigation.pdf}
    \caption{Autonomous aerial navigation from a language instruction: a UAV observation is grounded into a semantic-topometric graph; the instruction and graph condition an LLM planner that produces pose waypoints with velocity. The displayed route is a top-down projection, with attitude and altitude encoded in the waypoint state. Multi-drone planning is an extension with coordinated altitude separation. SR: success rate; SPL: success weighted by path length; NE: navigation error. Metric glyphs indicate evaluation criteria, not measured scores; scenes and graph geometry are schematic.}
    \label{fig:p4_s9_4_navigation}
\end{figure}


\subsection{Pivotal Task-3: Autonomous Aerial Navigation and Decision Making}
\label{subsec:p4_s9_4_task3}

The final stage closes the loop, converting perception into action. Traditional UAV geolocalization addresses basic aerial navigation by matching a query (\eg, a satellite image or textual description) to a database of historical flight paths or geographic points~\cite{zheng2020university,chu2024towards}. However, a critical distinction exists between such perception-oriented tasks and the overarching goal of autonomous navigation: the latter requires complex decision-making and control to bridge perception to action, drawing on the spatial-reasoning and planning paradigms of Chapter~\ref{sec:reasoning}. In practical scenarios, the objective is often for an aerial agent to autonomously locate a destination, particularly an unexplored target, through environmental exploration and, when available, collaboration with a map. This goal introduces significant challenges in generalization to unseen environments, including novel objects, unmapped regions, and interaction with the environment. For instance, recent benchmarks adapted from general Vision-Language Navigation (VLN) have begun to address this by requiring models to output a detailed flight plan with 6-DoF control signals~\cite{wang2024towards}. Figure~\ref{fig:p4_s9_4_navigation} follows the problem from the perception-to-action gap to foundation-model-driven VLN and multi-agent extensions.



\begin{table}[!ht]
\small
\centering
\caption{Comparison of recent aerial VLN datasets.}
\label{tab:aerial_datasets}
\begin{tabular}{@{}p{2.5cm}lp{2.5cm}p{2cm}p{3cm}p{3.5cm}@{}}
\toprule
\textbf{Dataset} & \textbf{Yr} & \textbf{Env.} & \textbf{Size} & \textbf{Modalities} & \textbf{Key Features} \\
\midrule
UAVA~\cite{albanis2020dronepose} & 2020 & Indoor simulated (MP3D) & 90 traj. & RGB, depth, flow, ego/exo views & Photorealistic UAV assist; vision/robotics tasks; gamified traj. \\
AVDN~\cite{fan2022aerial} & 2022 & Outdoor simulated (xView) & 3,064 traj. with dialog & RGB, dialog & Async human dialogs; attn. for interactive nav. \\
AerialVLN~\cite{liu2023aerialvln} & 2023 & Outdoor simulated (cities) & 8,446 traj. and 25,338 instructions & RGB, lang. instr. & Synthetic VLN; landmark nav in urban scenes. \\
CityNav~\cite{lee2024citynav} & 2024 & Real-world (UK cities) & 32,637 traj./instr. & 3D point clouds, RGB, geo-lang & Real 3D scans (4.65 km²); human demos; scalable VLN. \\
UAV-Need-Help~\cite{wang2024towards} & 2024 & Outdoor simulated (OpenUAV) & 12,149 traj. & RGB, lang. instr. & Continuous 6-DoF traj.; target-oriented. \\
OpenFly~\cite{gao2025openfly} & 2025 & Outdoor real/sim (18 scenes) & 100k traj. & 3D point clouds, RGB, lang. instr. & Large-scale; real-to-sim transfer. \\
AirV2X~\cite{gao2025airv2x} & 2025 & Outdoor simulated (CARLA+AirSim) & 6.73 hours of UAV-ground traj. & 3D point clouds, RGB & V2X collab; drone strategies; reduces occlusions. \\
\bottomrule
\end{tabular}
\end{table}


\textbf{From visual perception to language-guided planning.} Recent years have witnessed significant progress in UAV navigation, with a primary focus on visual perception~\cite{loquercio2018dronet,kang2019generalization,giusti2015machine,smolyanskiy2017toward,fan2020learn,bozcan2020air,majdik2017zurich} and collision avoidance~\cite{xu2025navrl,singla2019memory}. Building on this aerial-perception foundation, Vision-and-Language Navigation (VLN) for UAVs remains an emerging area of research. To address it, recent efforts have leveraged Large Multi-modal Models (LMMs) to harness general knowledge for navigation planning: STMR~\cite{gao2024aerial} builds upon AerialVLN~\cite{liu2023aerialvln} with a zero-shot LLM-based framework utilizing a Semantic-TopoMetric Representation for enhanced spatial reasoning. A parallel line reframes UAV autonomy as a full agentic loop rather than one-shot instruction following. AeroVerse~\cite{yao2024aeroverse} assembles a simulator, real-virtual pre-training corpora, and five downstream tasks, \ie, scene awareness, spatial reasoning, navigational exploration, task planning, and motion decision, to train aerospace embodied world models, and proposes SkyAgentX as a unified perception-reasoning-navigation-planning agent; AirVista-II~\cite{lin2025airvista} pursues zero-shot dynamic-scene semantic understanding via agent-based task scheduling and adaptive keyframe selection; and PRPSearcher~\cite{ji2025uavsearch} casts aerial visual object search over the CityAVOS benchmark as a perception-reasoning-planning problem solved by a multimodal LLM agent. Complementary systems such as NeLV~\cite{yuan2025nelv} operationalize the language-to-flight pipeline end to end, chaining an LLM parser, route and path planners, and a control platform across short-, medium-, and long-range missions. Beyond single-agent setups, some works have begun to explore multi-drone systems; for instance, Wu~\etal~\cite{wu2025aeroduo} deploy two drones flying at different altitudes to facilitate planning.

\textbf{Datasets and simulators.} As shown in Table~\ref{tab:aerial_datasets}, several foundational frameworks and datasets have been established. One of the early works, \ie, UAVA~\cite{albanis2020dronepose}, leverages the panorama images in the Matterport3D dataset~\cite{chang2017matterport3d} to render the indoor environment, while harnessing the camera viewpoints as physically available flight anchors to generate trajectories. AerialVLN~\cite{liu2023aerialvln} introduces a large-scale dataset with an effective baseline, and further infrastructure contributions include datasets such as CityNav~\cite{lee2024citynav}, AVDN~\cite{fan2022aerial}, and OpenFly~\cite{gao2025openfly}. Notably, UAV-Need-Help~\cite{wang2024towards} provides a realistic UAV dynamics simulator, \ie, OpenUAV, alongside its evaluation protocol; OpenUAV is further used in HaL-13k~\cite{wu2025aeroduo} to simulate two drones at different altitudes. Additional observation agents, such as road infrastructures and vehicles with depth sensing, have also been incorporated to enrich environmental understanding~\cite{gao2025airv2x}.

\textbf{Metrics.} Evaluation metrics for UAV aerial navigation, especially VLN, measure how well a model navigates a drone based on visual inputs and language instructions. Common metrics include Success Rate (SR), which checks if the UAV reaches the target location within a specified distance, and Success weighted by Path Length (SPL), which balances success with how efficiently the drone travels (shorter paths score higher). Path Length (PL) measures the total distance traveled, while Navigation Error (NE) calculates the distance between the final position and the goal. Together these metrics assess both accuracy (reaching the target) and efficiency (taking an optimal route), providing a clear way to compare model performance across datasets like UAVA, AerialVLN, and CityNav.


\subsection{Established Tools and Systems}
\label{subsec:p4-sec9-4-established-tools-and-systems}
Beyond the methods and benchmarks above, the aerial domain rests on a real-world layer of simulators, platforms, and infrastructure that turns research prototypes into deployable systems; we group it into three tiers.

\textbf{Simulation environments and benchmarks.} Because flight data is costly and safety-critical to collect, high-fidelity simulation underpins much of the field. The OpenUAV simulator released with UAV-Need-Help~\cite{wang2024towards} supplies realistic UAV dynamics and a continuous 6-DoF evaluation protocol, and is reused to simulate multi-altitude drone teams in HaL-13k~\cite{wu2025aeroduo}; OpenFly~\cite{gao2025openfly} targets large-scale real-to-sim transfer across diverse scenes; and AerialVLN~\cite{liu2023aerialvln} pairs a city-scale synthetic environment with a navigation baseline. Photorealistic indoor rendering for UAV tasks is provided through Matterport3D~\cite{chang2017matterport3d} in UAVA~\cite{albanis2020dronepose}, while V2X-style multi-agent simulation combines CARLA and Microsoft AirSim\footnote{\url{https://microsoft.github.io/AirSim/}} in AirV2X~\cite{gao2025airv2x}.

\textbf{Commercial platforms and vendors.} On the hardware side, consumer and enterprise drones from DJI\footnote{\url{https://www.dji.com/}} (whose Matrice series collects datasets such as UAV-Human) and autonomy-focused vendors such as Skydio\footnote{\url{https://www.skydio.com/}} dominate the UAV tier. In the stratospheric tier, solar-powered HAPS such as Airbus Zephyr~\cite{kothe2019flight} and AeroVironment/HAPSMobile's HAWK30~\cite{fladeland2019demonstrating} provide persistent regional coverage, complemented by small-satellite constellations built on CubeSat platforms~\cite{lemmer2017propulsion} and broader HAPS programs~\cite{gonzalo2018capabilities}.

\textbf{Connectivity and positioning infrastructure.} Aerial autonomy ultimately depends on shared positioning and communication backbones: the GNSS constellations introduced above (GPS~\cite{enge1994global}, GLONASS~\cite{ivanov1992glonass}, Galileo, and BeiDou~\cite{yang2019introduction}) supply global timing and localization, while space-based broadband relays such as Starlink~\cite{michel2022first} and orbital observation systems~\cite{cunningham2004national} increasingly serve as the connectivity layer for coordinated, real-time aerial operations.


\subsection{Challenges}
\label{subsec:p4-sec9-4-challenges}
Aerial autonomy is limited less by any single model than by the gaps between platforms, scales, and operating conditions. Four threads follow the pipeline in order: fusing heterogeneous platforms into a coherent situational model, perceiving small objects and complex activities reliably, navigating unseen and GPS-denied environments, and meeting the ethical and operational constraints under which any of this can be deployed.

\paragraph{Heterogeneous-Platform Fusion.}
A primary obstacle lies in integrating aerial agents with heterogeneous platforms for robust environmental perception. Coordinating a UAV with ground-based sensors or satellite systems to construct a cohesive situational-awareness model remains complex, particularly in dynamic or cluttered settings where the platforms differ in resolution, latency, and viewpoint (Table~\ref{tab:air-agents}). Reconciling these asynchronous, multi-resolution streams into a single consistent representation, without one modality's noise corrupting the fused estimate, is the standing difficulty, and it motivates the V2X-style cooperative-perception efforts that incorporate road infrastructure and ground vehicles as additional observers.

\paragraph{Small-Object and Complex-Activity Understanding.}
Aerial perception is defined by extremes of scale and density. The recognition of small objects, such as identifying debris in disaster zones or distinguishing wildlife in ecological surveys, demands high-resolution imaging and detectors robust to dense clustering, which often strain onboard computational resources. Understanding complex activities, such as discerning human behavior during search-and-rescue missions, compounds the problem by requiring nuanced spatio-temporal and contextual analysis that current systems struggle to deliver consistently. Closing this gap calls for detectors and activity models that remain accurate at extreme scale variation while fitting a UAV's tight compute and power budget.

\paragraph{Navigation in Unseen and GPS-denied Environments.}
Autonomous operation in unseen or GPS-denied environments, such as dense urban canyons or subterranean spaces, further complicates the action stage, as agents must rely on real-time mapping and obstacle avoidance without pre-existing data. The core requirement is generalization: a navigation policy must extend to novel objects, unmapped regions, and genuinely interactive scenarios, rather than replaying matched trajectories. This is precisely the gap that foundation-model-driven VLN and multi-altitude cooperative planning aim to narrow, but validating such behavior to a safety standard in the open world remains open.

\paragraph{Privacy, Airspace, and Operational Robustness.}
Beyond perception and navigation, several ethical and operational challenges persist. Data-privacy concerns arise when aerial agents capture high-fidelity multimedia data, potentially compromising individual rights without stringent safeguards, motivating privacy-preserving approaches such as federated learning. Airspace integration poses another hurdle, as regulatory frameworks struggle to accommodate the proliferation of aerial agents alongside traditional aviation, necessitating harmonized protocols to ensure safety. Energy efficiency remains a critical bottleneck, since prolonged missions demand lightweight yet powerful power sources and motivate renewable solutions such as solar-powered propulsion. Finally, operational robustness is challenged by environmental uncertainties, such as extreme weather or electromagnetic interference, that can disrupt system reliability. Advances in artificial intelligence and sensor fusion that address these constraints, rather than perception accuracy alone, are what will let aerial agents move from bounded demonstrations toward dependable, large-scale deployment.


\newpage

\section{Warehouse and Industrial Logistics}
\label{sec:p4_s9_5}



\subsection{Introduction}
\label{subsec:p4-sec9-5-introduction}

Warehouse and industrial logistics is one of the most mature real-world testbeds for the spatial capabilities surveyed in Part~\ref{part:3}: it demands perception-to-action pipelines that interpret complex, dynamic environments such as dense shelving systems populated with diverse stock keeping units (SKUs), pallets, forklifts, and human workers. 
These systems construct and maintain spatial memory representations (semantic maps, object and pose states) to support downstream planning and manipulation for tasks including inventory counting, order picking, pallet handling, and coordinated fleet operations. 
Perception and reasoning themselves are treated in Chapter~\ref{chapter:3:sec:6} and Chapter~\ref{sec:reasoning}; logistics adds throughput targets, engineered but crowded layouts, and people sharing the floor with fleets of robots. 
Performance is judged in terms of throughput, order accuracy, safety, and operational cost, especially under high SKU diversity, visual clutter, elongated aisles, and human-robot collaboration. 
Dense shelving and manipulation in clutter remain the hardest cases, as the Amazon Picking Challenge made plain~\cite{Correll2016AmazonChallenge}, and they demand that perception, pose estimation, mapping, and multi-agent planning be integrated rather than tuned in isolation.

\begin{figure}[!ht]
    \centering
    \includegraphics[width=0.98\linewidth]{figures/part4/section9/warehouse_overview.pdf}
    \caption{Spatial intelligence for warehouse and industrial logistics: RGB-D, LiDAR, and RFID observations support five complementary capabilities for picking and pallet transport, and execution outcomes update maps and plans. Scene images and diagrams are illustrative.}
    \label{fig:p4_s9_5_overview}
\end{figure}

\subsection{System Overview}
\label{subsec:p4-sec9-5-system-overview}

The complete warehouse automation pipeline, sketched in Figure~\ref{fig:p4_s9_5_overview}, integrates five complementary capabilities into a closed loop. 
\textbf{Perception} converts raw multi-sensor streams into semantic object information (what is present, where, and in what pose); 
\textbf{localization and mapping} grounds that information in a metric spatial map and keeps it consistent over time; 
\textbf{reasoning and planning} consume these representations to allocate tasks and generate optimized fleet schedules; 
\textbf{manipulation and control} physically execute the resulting picking and transport actions; 
and a \textbf{human-robot collaboration and safety} layer adapts behavior dynamically in mixed human-robot environments. 
Integration is achieved via middleware such as ROS\,2 and industrial protocols like Open Platform Communications Unified Architecture (OPC-UA), while digital twin platforms such as NVIDIA Omniverse\footnote{https://www.nvidia.com/en-gb/omniverse/} support simulation and validation before deployment. 
The five capabilities are the pivotal tasks below, taken in pipeline order.



\subsection{Pivotal Task-1: Perception and Recognition}
\label{subsec:p4-sec9-5-pivotal-task-1-perception-and-recognition}
Perception constitutes the sensory foundation of the pipeline: it transforms raw sensor data into actionable knowledge, reusing the detection, segmentation, and pose-estimation machinery of Chapter~\ref{chapter:3:sec:6} but specializing it to the warehouse Operational Design Domain. 
Core capabilities include detecting and classifying objects (packages, pallets, tools), estimating their 3D poses, and tracking their motion and state over time. 
Unlike outdoor environments, warehouses impose distinct perceptual challenges: dense clutter, heavy occlusion, visually similar items, reflective packaging, narrow aisles, and continually changing layouts, which together make robust perception a central and ongoing research focus \cite{Tubis2023IntelligentWarehouse,redmon2016yolo}. 
Formally, perception is a structured mapping from a synchronized multimodal stream $X_t = \{I_t^{\text{RGB}}, I_t^{\text{Depth}}, P_t^{\text{LiDAR}}, R_t^{\text{RFID}}\}$ (together with calibration parameters and a prior map $M$) to a scene graph $G_t = \{o_i, c_i, \mathbf{b}_i, \mathbf{p}_i, \Sigma_i\}$, in which each detected object $o_i$ carries a class label $c_i$, a bounding box or segmentation mask $\mathbf{b}_i$, an estimated pose $\mathbf{p}_i$, and an uncertainty measure $\Sigma_i$ that grounds subsequent manipulation, navigation, and decision making. 
Figure~\ref{fig:p4_s9_5_perception} sketches this representation stack.

\begin{figure}[!ht]
    \centering
    \includegraphics[width=0.98\linewidth]{figures/part4/section9/warehouse_perception.pdf}
    \caption{The warehouse perception stack: 2D detection and segmentation, 6D object pose $T_i$, and a structured scene graph whose nodes carry identity, class, box or mask, pose, and uncertainty $\Sigma_i$. RFID can support identity association; masks and axes are schematic, not calibrated estimates.}
    \label{fig:p4_s9_5_perception}
\end{figure}

\textbf{Detection and recognition} form the base layer. 
Convolutional and transformer-based architectures, \ie, YOLO \cite{redmon2016yolo}, RetinaNet \cite{lin2017focal}, and the transformer-based DETR with deformable attention \cite{Zhu2021DeformableDETR}, have steadily improved accuracy and efficiency under clutter, and are scored on COCO \cite{lin2014microsoft}, whose mean Average Precision (mAP) broken down by object size is particularly informative for the small-SKU and barcode-label cases that dominate warehouse failure modes. 
\textbf{Pose estimation} lifts these 2D detections into the metric space that manipulation requires: PoseCNN \cite{xiang2017posecnn}, DenseFusion \cite{Wang2019DenseFusion}, and CosyPose \cite{Labbe2020CosyPose} fuse visual and geometric cues to infer six-degree-of-freedom (6D) orientation and position even under occlusion, and are evaluated on YCB-Video \cite{xiang2017posecnn} for textured household and industrial objects and on T-LESS \cite{Hodan2017TLESS} for textureless, industry-relevant parts, with Average Distance (ADD) the standard accuracy measure and MOTA/MOTP reporting tracking quality in dynamic scenes. 
The Amazon Picking Challenge \cite{Correll2016AmazonChallenge} couples perception with manipulation in real warehouse scenarios, stressing generalization across novel object sets and thereby anticipating the manipulation task of §\ref{subsec:p4_s9_5_manip} below. 
Representative benchmarks across the section are consolidated in Table~\ref{tab:p4_s9_5_benchmarks}. In industrial practice, perception pipelines are increasingly standardized and modularized; the deployed software stacks that realize them are surveyed in the Established Tools and Systems subsection.

\subsection{Pivotal Task-2: Localization and Mapping}
\label{subsec:p4-sec9-5-pivotal-task-2-localization-and-mapping}
If perception answers \emph{what is around the robot}, localization and mapping answer \emph{where the robot is} and \emph{what the persistent structure of the warehouse is}, forming the spatial backbone on which fleet coordination rests and drawing on the same scene-construction techniques as Chapter~\ref{chapter:3:sec:6}. 
Automated guided vehicles (AGVs) and autonomous mobile robots (AMRs) require centimeter-level accuracy to navigate narrow aisles, dock precisely at shelves, and operate safely alongside workers. 
The problem is conceptually the SLAM problem of autonomous driving, yet it faces warehouse-specific constraints: highly repetitive structures (rows of identical shelves), dynamic obstacles (workers, pallets, forklifts), and low-light or GPS-denied indoor conditions. 
Formally, a sequence of odometry and sensor measurements $Z_{1:t} = \{\text{IMU}, \text{LiDAR}, \text{cameras}, \text{wheel encoders}, \text{UWB beacons}\}$, optionally supplemented by a prior map $M_0$ or fiducial anchors, is fused into a pose graph $\mathcal{X}_t = \{x_1, \ldots, x_t\}$ that estimates the trajectory and an evolving semantic map $M_t$ fusing geometric structure with storage zones and restricted areas. 
Beyond navigation, this map underpins the digital twins (spatial models synchronized with the facility whose simulations feed back into its operation; §\ref{sec:p4_s11_1}) that increasingly govern modern facilities \cite{Cadena2016SLAMSurvey}. Figure~\ref{fig:p4_s9_5_localization} casts the task as a multi-sensor fusion problem.

\begin{figure}
    \centering
    \includegraphics[width=0.98\linewidth]{figures/part4/section9/warehouse_localization.pdf}
    \caption{Localization and mapping in the warehouse as a factor-graph fusion problem: IMU, LiDAR, camera, wheel-encoder, and UWB measurements constrain pose states, optionally with a prior map or beacon anchors, producing poses over warehouse geometry, storage zones, and restricted areas. Repetitive shelving can cause false place matches, so only verified loop candidates add constraints.}
    \label{fig:p4_s9_5_localization}
\end{figure}

A family of \textbf{SLAM algorithms} has been tailored to these conditions. ORB-SLAM3 \cite{Campos2021ORBSLAM3} extends visual SLAM with inertial and multi-camera support for accurate visual-inertial localization; LIO-SAM \cite{Shan2020LIOSAM} provides real-time LiDAR--IMU fusion through factor-graph optimization that stays robust under fast motion and partial occlusion; and RTAB-Map \cite{Labbe2019RTABMap} supports online RGB-D mapping and loop-closure detection with memory management suited to long-term operation, all expressible within the unified factor-graph framework of GTSAM \cite{Dellaert2012GTSAM}. 
\textbf{Multi-modal fusion} is the common thread that holds accuracy together when any single modality degrades, and is also the dominant emerging research direction, alongside learning-based estimation (neural depth prediction, differentiable bundle adjustment, and learned place-recognition embeddings) and the migration of SLAM computation to the edge for fleet-scale latency. 
These capabilities are evaluated on KITTI \cite{Geiger2013KITTI} for outdoor visual and LiDAR odometry and loop closure, on the indoor visual-inertial EuRoC MAV dataset \cite{Burri2016EuRoC} whose challenging motion and illumination mirror warehouse conditions, and on the high-fidelity Replica reconstructions \cite{straub2019replica} for dense mapping and semantic labeling, with Absolute and Relative Pose Error (APE/RPE) and loop-closure precision/recall the standard metrics and, for warehouse use, semantic-zone IoU and map completeness (Table~\ref{tab:p4_s9_5_benchmarks}).


\begin{figure}[!ht]
    \centering
    \includegraphics[width=0.98\linewidth]{figures/part4/section9/warehouse_planning.pdf}
    \caption{Two families for warehouse reasoning and planning over shared task, fleet, and map inputs: combinatorial optimization (assignment and routing) and learning-based coordination (per-agent policies), optionally combined in hierarchical or hybrid designs. Allocations, paths, and schedules are illustrative.}
    \label{fig:p4_s9_5_planning}
\end{figure}



\subsection{Pivotal Task-3: Reasoning and Planning}
\label{subsec:p4-sec9-5-pivotal-task-3-reasoning-and-planning}
Reasoning and planning form the decision-making core of warehouse automation, transforming the perceptual and localization outputs into task assignments and motion plans, and instantiating the spatial-reasoning and planning paradigms of Chapter~\ref{sec:reasoning}. 
This layer determines which robot executes which task, how tasks are sequenced, and how routes and schedules are optimized for safe, efficient, coordinated operation. 
In contrast to autonomous driving, where planning typically focuses on individual vehicles, warehouse reasoning must address multi-agent coordination across fleets of dozens or hundreds of robots in constrained indoor space \cite{Veloso2015CMUCoBot}, maximizing throughput and minimizing latency while maintaining safety under dynamic obstacles, human interactions, and stochastic task arrivals. 
Formally, the perceived states $G_t$, a task list $\mathcal{T}$, the fleet state $F_t$ (positions, capabilities, battery levels), and the warehouse map $M_t$ are mapped to a task-allocation matrix $A$, optimized motion paths $\mathcal{R}$, and temporal schedules $\pi$ that jointly specify when and how tasks execute. Figure~\ref{fig:p4_s9_5_planning} contrasts the two solution families.


The first family is \textbf{combinatorial optimization}. Traditional methods (the Hungarian algorithm, branch-and-bound, and mixed-integer linear programming) remain widely used for task allocation and scheduling in small- to medium-scale deployments, while Vehicle Routing Problem (VRP) formulations \cite{Cordeau2007VRP} and their multi-depot and time-window variants are central to modeling fleet movement. The second family is \textbf{learning-based coordination}: reinforcement learning (RL) and multi-agent RL (MARL) handle dynamic environments and real-time uncertainty \cite{Sartoretti2019PRIMAL,samvelyan2019smac}, with toolkits such as Google OR-Tools \cite{Perron2019ORTools} and PyMARL enabling rapid experimentation across both families. 
Emerging research couples the two: hierarchical and hybrid frameworks combine symbolic reasoning with learned components, graph neural networks and differentiable optimization accelerate multi-agent coordination and transfer policies across layouts, and digital-twin integration enables simulation-informed, predictive re-planning, increasingly under multi-objective criteria that balance throughput, energy, and human safety. 
Planning systems are benchmarked on Solomon's VRP suite \cite{Solomon1987VRP} for static routing and scheduling and on the Open Motion Planning Library (OMPL) \cite{Sucan2012OMPL} for continuous-space motion planning, complemented by synthetic multi-robot testbeds; common metrics are throughput (items/hour), task-completion latency, total travel distance or time, collision rate, and re-planning latency (Table~\ref{tab:p4_s9_5_benchmarks}).

\subsection{Pivotal Task-4: Manipulation and Control}
\label{subsec:p4_s9_5_manip}
Manipulation constitutes the physical execution layer, enabling robots to interact directly with objects, containers, and infrastructure such as bins, shelves, and conveyors. 
In industrial logistics this typically involves robotic arms performing grasping, placing, sorting, or palletizing, as well as AGVs docking to pallets or workstations. Manipulation bridges high-level symbolic or geometric plans with continuous motor control, and therefore demands both precision and robustness to the inherent variability in packaging, geometry, and material properties \cite{Levine2016DeepGrasp}; the grasping and control techniques here connect directly to the robotic-manipulation treatment of §\ref{sec:p4_s9_1} and the planning paradigms of Chapter~\ref{sec:reasoning}. 
Formally, estimated object poses $\mathbf{p}_i$ and task descriptions (grasp targets, placement goals, trajectory constraints) are mapped to low-level motor commands $\mathbf{u}_t$ that ensure stable grasping, precise placement, and compliant, collision-free motion. Figure~\ref{fig:p4_s9_5_manipulation} contrasts modular grasp-and-plan systems with learned visuomotor control.

\begin{figure}
    \centering
    \includegraphics[width=0.98\linewidth]{figures/part4/section9/warehouse_manipulation.pdf}
    \caption{Warehouse manipulation through (a) a modular route, where grasp candidates scored by analytic metrics and learned quality (Dex-Net) feed sampling-based planning (RRT*) and trajectory tracking, and (b) learned visuomotor control from camera observations. Action spaces vary across systems, and the routes can be combined.}
    \label{fig:p4_s9_5_manipulation}
\end{figure}

\textbf{Grasp synthesis} anchors the methodology: Dex-Net combines synthetic grasp data and analytic robustness metrics to train a grasp-quality predictor for candidate selection \cite{Mahler2017DexNet}, while learned visual control can predict grasp success from camera images and candidate task-space motions, then select actions for adaptive manipulation \cite{Levine2016DeepGrasp}. 
\textbf{Motion planning} then realizes the chosen grasp: sampling-based algorithms such as RRT* \cite{Karaman2011RRTStar} remain fundamental for generating collision-free trajectories in high-dimensional configuration spaces, exposed through frameworks like MoveIt \cite{Sucan2010MoveIt}. 
These capabilities are benchmarked on the Amazon Picking Challenge \cite{Correll2016AmazonChallenge}, which established a standardized testbed for picking under clutter and uncertainty, and on the YCB Object and Model Set \cite{Calli2015YCB}, a diverse collection of everyday items with physical models for grasping and manipulation; grasp success rate, cycle time per pick, placement accuracy, and task-completion rate are the headline metrics (Table~\ref{tab:p4_s9_5_benchmarks}). 
Hybrid control frameworks that combine analytical models with data-driven adaptation are emerging as the practical route to generalizing across variable geometries and grasp contexts, complemented by tactile sensing and force control for slip detection and contact regulation.

\subsection{Pivotal Task-5: Human-Robot Collaboration and Safety}
\label{subsec:p4-sec9-5-pivotal-task-5-human-robot-collaboration-and-safety}
Even in highly automated warehouses, human workers remain indispensable for complex, dexterous, or exception-handling tasks that require contextual understanding and flexible decision making, so the final stage of the pipeline is safe and efficient human-robot collaboration (HRC). 
HRC systems must balance productivity with stringent safety guarantees, enabling mixed teams (robots handling repetitive or heavy-load work, humans managing fine-grained, judgment-based activities) to share workspaces without physical or operational interference \cite{Lasota2015HRC}. 
The task maps multimodal human-detection and intent signals (pose, gesture, gaze, and trajectory data from vision, LiDAR, wearables, or motion capture) to adaptive robot behaviors such as dynamic speed modulation, path re-routing, stop-and-wait actions, or task reallocation in response to predicted human movement. Figure~\ref{fig:p4_s9_5_hrc} sketches this predict-then-adapt loop.

\begin{figure}[!ht]
    \centering
    \includegraphics[width=0.98\linewidth]{figures/part4/section9/warehouse_human_robot_collaboration.pdf}
    \caption{The human--robot collaboration loop: multimodal human sensing and intent recognition inform proximity-aware planning, which selects slowing, rerouting, or stopping. Clearance regions and routes are illustrative, not certified safety zones.}
    \label{fig:p4_s9_5_hrc}
\end{figure}

\textbf{Intent recognition} lets robots anticipate human motion: sequence models based on hidden Markov models or recurrent networks predict trajectories and proactively adjust behavior \cite{Mainprice2016Intent}, while \textbf{proxemics-aware planning} maintains socially acceptable distances and re-routes trajectories under human proximity \cite{dragan2013legibility}, with multimodal sensing (vision-based detection fused with wearable or RFID localization) improving reliability. 
These principles are embodied in force- and speed-limited collaborative robots and safety-rated AMRs (see Established Tools and Systems), governed by international standards such as ISO 10218 and ISO/TS 15066. 
Evaluation combines quantitative safety metrics with qualitative interaction assessment: human-motion corpora such as Human3.6M~\cite{IonescuPOS14} supply sequences for intent prediction, and key indicators include intervention frequency (safety-triggered stops per hour), near-miss rate, human-detection accuracy, and compliance with ISO-defined safety zones.

\begin{table}[!th]
\centering
\fontsize{8}{10}\selectfont
\setlength{\tabcolsep}{1.0mm}
\caption{Representative benchmarks for spatial intelligence in warehouse and industrial logistics, spanning the five pivotal tasks.}
\label{tab:p4_s9_5_benchmarks}
\begin{tabular}{L{0.21\textwidth} L{0.21\textwidth} L{0.31\textwidth} L{0.18\textwidth}}
\toprule
\textbf{Pivotal task} & \textbf{Benchmark} & \textbf{Representative method} & \textbf{Headline metric} \\
\midrule
Perception (detection) & COCO~\cite{lin2014microsoft} & YOLO~\cite{redmon2016yolo} / DETR~\cite{Zhu2021DeformableDETR} & mAP \\
Perception (6D pose) & YCB-Video~\cite{xiang2017posecnn} & DenseFusion~\cite{Wang2019DenseFusion} & ADD \\
Perception (textureless pose) & T-LESS~\cite{Hodan2017TLESS} & CosyPose~\cite{Labbe2020CosyPose} & ADD / VSD \\
Localization \& mapping & KITTI~\cite{Geiger2013KITTI} & Visual / LiDAR SLAM & APE / RPE \\
Localization \& mapping & EuRoC MAV~\cite{Burri2016EuRoC} & ORB-SLAM3~\cite{Campos2021ORBSLAM3} (visual-inertial) & APE \\
Mapping (dense) & Replica~\cite{straub2019replica} & RTAB-Map~\cite{Labbe2019RTABMap} & IoU / completeness \\
Reasoning \& planning & Solomon VRP~\cite{Solomon1987VRP} & OR-Tools~\cite{Perron2019ORTools} & total cost / distance \\
Reasoning \& planning & OMPL~\cite{Sucan2012OMPL} & Sampling-based planners & planning time / success \\
Manipulation & Amazon Picking~\cite{Correll2016AmazonChallenge} & Integrated perception--grasping team pipelines & pick success / time \\
Manipulation & YCB Object Set~\cite{Calli2015YCB} & Dex-Net~\cite{Mahler2017DexNet} & grasp success rate \\
\bottomrule
\end{tabular}
\end{table}

\subsection{Established Tools and Systems}
\label{subsec:p4-sec9-5-established-tools-and-systems}
The pivotal tasks above are realized in production through a layered ecosystem of simulators, middleware, optimization services, and robot hardware that together constitute the real-world delta over the methods of Part~\ref{part:3}.

\textbf{Simulation and middleware} form the integration substrate. Digital-twin and physics platforms such as NVIDIA Omniverse Isaac Sim \cite{NVIDIA2021Isaac} and Gazebo provide realistic environments for developing and validating perception, localization, and planning pipelines before deployment, while perception SDKs such as NVIDIA Isaac, the Intel RealSense SDK, and Zebra RFID/barcode platforms  supply end-to-end sensing support in real facilities. Interoperability across heterogeneous hardware is provided by ROS~\cite{Quigley2009ROSIndustrial}, ROS\,2, and ROS-Industrial, with commercial localization stacks such as SLAMcore\footnote{\url{https://www.slamcore.com/}} and BlueBotics ANT\footnote{\url{https://bluebotics.com/}} delivering optimized multi-sensor fusion, map management, and fleet coordination for AGVs and AMRs.

\textbf{Warehouse management and fleet orchestration} sit at the enterprise layer. Large-scale optimization services such as SAP Extended Warehouse Management\footnote{\url{https://www.sap.com/products/scm/extended-warehouse-management.html}}, Blue Yonder WMS\footnote{\url{https://blueyonder.com/}}, and Manhattan Associates\footnote{\url{https://www.manh.com/}} integrate task allocation and scheduling with enterprise resource planning, while open toolkits such as Google OR-Tools \cite{Perron2019ORTools} and PyMARL support the underlying combinatorial and learning-based planners.

\textbf{Robot hardware and manipulation} close the loop physically. High-throughput pick-and-place arms from ABB\footnote{\url{https://new.abb.com/products/robotics}}, FANUC\footnote{\url{https://www.fanuc.com/}}, and KUKA\footnote{\url{https://www.kuka.com/}} are widely deployed and integrated through ROS-based middleware, mobile manipulators such as those from Fetch Robotics\footnote{\url{https://fetchrobotics.com/}} unify navigation and grasping, and motion-planning frameworks such as MoveIt \cite{Sucan2010MoveIt} and NVIDIA cuMotion provide inverse kinematics and trajectory optimization. For shared workspaces, collaborative robots from Universal Robots\footnote{\url{https://www.universal-robots.com/}} and AMRs from Locus Robotics\footnote{\url{https://locusrobotics.com/}} provide the safety-rated platforms on which the HRC layer is built.






\subsection{Challenges}
\label{subsec:p4-sec9-5-challenges}
Deployed warehouse systems work because they fix the SKU sets, layouts, and operating conditions they handle, and each stage of the pipeline has its own obstacle to relaxing those limits: small-object and cluttered perception, perceptual aliasing in repetitive aisles, multi-agent scalability under uncertainty, manipulation of variable objects, and the safety--throughput trade-off in human--robot collaboration.

\paragraph{Cluttered and Small-Object Perception.}
Barcode labels and compact SKU packages often occupy only a few pixels in high-resolution images, leading to frequent misdetections, while occlusion and clutter compound the difficulty as items overlap within bins and shelves. 
A further persistent issue is domain shift: models trained on synthetic or generic data degrade under real-world variation in lighting, reflective surfaces, and packaging. 
Research responses include domain adaptation and self-supervised learning for generalization \cite{Tzeng2017DomainAdversarial}, active perception that deliberately adjusts viewpoints or sensing to resolve occlusion \cite{Bajcsy2018ActivePerception}, and continual learning that lets perception adapt over time without exhaustive retraining, increasingly combined with multi-sensor fusion and neuro-symbolic representations for higher-level causal and semantic reasoning.

\paragraph{Perceptual Aliasing in Repetitive Environments.}
Localization and mapping are degraded by dynamic obstacles (humans, forklifts) that generate transient or misleading landmarks corrupting loop-closure hypotheses \cite{Cadena2016SLAMSurvey}, and above all by perceptual aliasing: repetitive shelving and low-texture surfaces produce ambiguous place recognition and false data associations in feature-based SLAM \cite{Campos2021ORBSLAM3,Labbe2019RTABMap}. 
High-speed AGV and AMR operation further demands low-latency, sub-centimeter pose updates that sensor delays, occlusion, and monocular scale drift can readily violate. Robust multi-modal fusion, learning-based estimation, and tighter integration of SLAM with semantic perception and digital twins are the leading mitigations.

\paragraph{Scalability and Uncertainty in Multi-agent Planning.}
The combinatorial complexity of multi-robot coordination grows exponentially with the number of agents and tasks, rendering exact optimization intractable for large fleets, while real-world uncertainty (new order arrivals, blocked aisles, unpredictable human movement) necessitates adaptive re-planning and decentralized decision making under partial observability \cite{Veloso2015CMUCoBot,Sartoretti2019PRIMAL}. 
Hierarchical and hybrid planners, graph-neural and differentiable coordination, and simulation-informed multi-objective optimization are the active directions toward scalable, resilient fleet behavior.

\paragraph{Manipulation of Variable and Deformable Objects.}
Many SKUs exhibit deformable, reflective, or transparent surfaces that interfere with perception and grasp stability \cite{Calli2015YCB,Mahler2017DexNet}, and deformable handling remains especially difficult because grasp dynamics and contact points resist rigid-body modeling \cite{Levine2016DeepGrasp}. 
Tactile sensing and force control improve slip detection and grasp robustness but face cost and calibration barriers in industrial deployment, motivating hybrid control that combines analytical models with data-driven adaptation for generalization across geometries and grasp contexts.

\paragraph{Safety--Throughput Trade-off in Collaboration.}
Robust human detection under occlusion, variable lighting, and cluttered aisles continues to hinder reliable intent estimation \cite{dragan2013legibility,Mainprice2016Intent}, and the trade-off between throughput and safety remains unresolved: conservative margins reduce efficiency while aggressive planning risks unsafe interaction \cite{Lasota2015HRC}. 
Emerging work emphasizes predictive and explainable safety models that justify decisions in human-interpretable terms, and adaptive safety frameworks that fuse learned intent models with certified control policies for dynamic risk assessment and shared autonomy on the warehouse floor.
